\mysection{\textit{OnlineGN}: Симулатор за обобщеномрежови модели}

\hspace*{2.5em}В настоящата глава се представя подробно разработеният в рамките на дисертационния труд нов симулатор за ОМ, като се разглеждат неговите изисквания, архитектурни решения, проектен дизайн и ключови функционалности. Системата представлява самостоятелен и изцяло нов софтуерен компонент, който към момента не е интегриран с други съществуващи решения. Предвижда се обаче в бъдещи етапи да бъде реализирана възможност за работа с файлови формати, съдържащи ОМ, генерирани от други симулатори, с цел повишаване на съвместимостта и улесняване на миграцията на съществуващи модели. Изходният код на приложението е публикуван в \cite{phdThesisOnlineGNRepo}.

Всички функционалности на софтуерната система са представени нагледно в подготвена видео демонстрация \cite{OnlineGN-video}, която илюстрира практическия цикъл на работа с разработения симулатор.

\subsection{Изисквания и цели на системата}

\hspace*{2.5em}Този раздел е посветен на изпълнението на \ref{task:requirements}, като дефинира критерии, на които трябва да отговаря предложеното уеб базиранo решение. Изискванията към системата се разделят на два основни типа: функционални и нефункционални.

\subsubsection{Функционални изисквания}
\hspace*{2.5em}Функционалните изисквания описват конкретните възможности и поведения, които системата трябва да предоставя на потребителя с цел изграждане, управление и симулация на ОМ модели.

\textbf{Системата трябва да:}
\begin{itemize}[itemsep=0.5em, topsep=0pt]
	\item визуализира изчертана в двумерното пространство мрежа с нейните преходи, позиции и коректните връзки между тях;
	
    \item извършва симулация на функционирането на ОМ \textbf{чрез прилагане на алгоритъм Б} (дефиниран в \textbf{\ref{tokenMove}}), регламентиращ активирането на преходите и движението на ядрата между позициите. Симулацията обхваща промяна на характеристиките, разделяне и сливане на ядра, като се спазва правилният ред на обработка на преходи/състояния/ядра. Това трябва да е възможно по два различни начина:
    \begin{itemize}
        \item \textbf{Режим за проследяване} -- изпълнение на една симулационна стъпка при поискване;
        \item \textbf{Непрекъснато изпълнение} -- автоматично изпълнение на всички симулационни стъпки до достигане на крайно състояние (или зададен критерий за спиране).
    \end{itemize}

	
	\item предоставя рестартиране на симулацията, при което всяко ядро се връща в първоначалната си позиция и състояние;
	
	\item позволява промяна на координатите на изчертаните преходи, позиции и връзките между тях, включително:
	\begin{itemize}
		\item промяна на координатите на позиции;
		\item промяна на координатите и дължините на преходите;
		\item промяна на координатите на точките, през които минава връзка между преход и позиция, както и добавяне/премахване на точки.
	\end{itemize}
	
	\item визуализира и позволява промяна на характеристиките на преходите и позициите, като това се отразява върху симулацията;
	
	\item предоставя подробна информация за характеристиките и местоположението на ядрата през всяка стъпка от симулацията;
	
	\item позволява добавяне на нови потребителски функции (на \textit{JavaScript}), които се изпълняват при преходи и позиции, и имат достъп до всички ядра, включително текущото;
	
	\item позволява динамична промяна на скоростта на симулацията;
	
	\item предоставя регистър, проследяващ последователността на операциите в рамките на една времева стъпка. Регистърът трябва да отразява реда на обработка на компонентите (преходи, позиции и ядра) в съответствие с техните приоритети и капацитети, като изяснява логиката на задействане на преходите и движението на ядрата;
	
	\item зарежда мрежи от \textit{JSON} файл в дефинирания от системата формат;
	
	\item записва мрежи в \textit{JSON} файл в дефинирания формат;
	
	\item поддържа зареждане на съхранени мрежи от база данни;
	
	\item съхранява променени мрежи и предоставя хипервръзка за достъп.
	
	\item генерира представяне на мрежата в \textit{TeX} формат на база предварително зададени параметри;
	
	\item изтрива автоматично мрежи от сървъра, които не са били достъпвани за по-дълъг период от време от предварително дефинирания.
\end{itemize}

\vspace{0.25cm}

\subsubsection{Нефункционални изисквания}
\label{sec:non-functional-requirements}
\hspace*{2.5em}Нефункционалните изисквания дефинират качествените характеристики на системата, които не са пряко свързани с конкретна функционалност, но оказват съществено влияние върху нейната надеждност, ефективност, сигурност и използваемост.

\textbf{Системата трябва да:}
\begin{itemize}[itemsep=0.5em, topsep=0pt]
    \item осигурява \textbf{висока производителност} – времето за отговор на всяка заявка да не надвишава 2 секунди при ефективно използване на ресурсите;
    \item осигурява \textbf{висока надеждност} – времето за възстановяване след отказ да не надвишава 30 минути;
    \item гарантира \textbf{сигурност} – използване на защитени протоколи за комуникация и криптиране на данните;
    \item осигурява \textbf{висока използваемост} – потребителският интерфейс да бъде интуитивен, последователен и лесен за работа;
    \item бъде \textbf{мащабируема} – архитектурата позволява лесно добавяне на нови функционални модули;
    \item \textbf{не събира лични данни} и да бъде в съответствие с изискванията на \textit{GDPR}.
\end{itemize}
\vspace{0.4cm}
\textbf{Обосновка на избраните параметри за време за отговор и възстановяване.} \\
Прагът за време за отговор \textbf{до 2 секунди} е дефиниран като практична горна граница за интерактивна работа при \emph{леко натоварване} (например ограничен брой едновременно активни потребители и типични по обем входни данни). Целта е да се постигне баланс между използваемост и техническа реализуемост, така че системата да се възприема като достатъчно отзивчива в основните сценарии на работа. Фундаментални изследвания в областта на \textit{HCI/UX} установяват, че около 1 секунда е ориентировъчна граница за запазване на „потока“ на мисълта, докато по-дългите забавяния изискват допълнителна обратна връзка и започват да влияят негативно върху субективната оценка за отзивчивост \cite{Nielsen1993RT}. Емпирични данни от контролирано изследване на латентност при търсене допълнително показват, че при увеличаване на забавянето до няколко секунди потребителските предпочитания се изместват осезаемо към по-бързи системи, като се наблюдава ясно влошаване при преминаване над приблизително 3 секунди \cite{Brutlag2008Latency}. На тази основа \textbf{2 секунди} се приема като консервативна, измерима и проверима целева стойност за горна граница в леки работни сценарии. Формулировката „ефективно използване на ресурсите“ следва терминологията на \textit{ISO/IEC~25010}, където \textit{performance efficiency} се разглежда като производителност спрямо използваните ресурси при зададени условия \cite{ISOIEC25010_2011,Arc42ISO25010}.

Прагът за възстановяване след отказ \textbf{до 30 минути} е зададен като цел от тип \textit{RTO (Recovery Time Objective)} -- максимално допустимият период, през който компонентите на системата могат да се намират във фаза на възстановяване преди да настъпи неблагоприятно въздействие върху целите/процесите \cite{NIST80034r1,NISTRTOGlossary}. В контекста на разработваната система (изследователски софтуер/симулатор при сравнително леко натоварване) стойността \textbf{30 минути} представлява разумна цел, която осигурява предвидимост на възстановяването, без да налага внедряване на скъпоструващи архитектури за непрекъсваема работа. Допълнително, прагът е съвместим с широко използвана цел за наличност от порядъка на 99.9\% за 30-дневен период, при която допустимият престой е приблизително 43,2 минути; така \textit{RTO} = 30 минути се вписва в приемлив „бюджет за престой“ и оставя резерв за диагностика и планирани дейности \cite{AtlassianSLO,BetterStackAvailability}.

\vspace{0.25cm}

\subsubsection{Ограничения и предпоставки на средата}

\hspace*{2.5em}Тъй като системата е внедрена, използвайки безплатните облачни услуги на \textbf{\textit{Render}} (за сървъра) и \textbf{\textit{Atlas}} (за базата данни), тя трябва да се съобразява с ограничените ресурси, предоставени от тези услуги.

\vspace{0.4cm}

\subsection{Архитектура на системата} 
\hspace*{2.5em}В този раздел се представя проектирането на архитектурата на системата \textit{OnlineGN}, което съответства на изпълнението на първа част от \ref{task:architecture}. Разглежда се архитектурната организация, като се описват основните компоненти, структурното разделение на модулите и начините на взаимодействие между тях.

\subsubsection{Обща архитектура и сценарии на взаимодействие}
\hspace*{2.5em}
Както е илюстрирано на \textbf{Фигура~\ref{fig:arch-layers-overview}}, системата е реализирана като трислойна архитектура, включваща слой за данни (база данни), сървърен слой (\textit{API}/услуги) и клиентски слой (потребителско приложение), при което клиентът поема основната част от бизнес логиката, а сървърът предоставя спомагателни услуги.

\begin{figure}[H]
	\centering
    \captionsetup{width=0.8\textwidth}
	
\includegraphics[width=0.7\linewidth]{"3. onlineGn/images/arch-layers-overview.png"}
	\caption{Диаграма, представяща трислойната архитектура на системата (клиент, сървър, БД) и потока на взаимодействие между основните компоненти.}
	\label{fig:arch-layers-overview}
\end{figure}

В контекста на тази архитектура, всеки от трите компонента има следните специфични отговорности:

\paragraph{Слой на базата от данни.}
\mbox{}\\
Този слой съхранява всички данни на приложението (мрежи, състояния и др.). Той позволява прочитане, извличане и управление на данните. Основната му роля е да отговаря за консистентността и цялостта на информацията.

\paragraph{Сървърен слой.} 
\mbox{}\\
Този слой служи като посредник (\textit{API}) между клиентското приложение и базата от данни. Освен това, той изпълнява задачи по поддръжка, като например периодично изтриване на неизползвани или остарели записи от базата данни.

\paragraph{Клиентски слой.} 
\mbox{}\\
Клиентският слой представлява потребителския интерфейс (\textit{UI/UX}) на приложението. В него се извършва цялата бизнес логика, свързана с визуализацията на мрежата, симулацията, модифицирането на елементи и други интерактивни действия.

\paragraph{Сценарии на взаимодействие.}
\mbox{}\\
Следните сценарии описват комуникационния поток между слоевете при основните потребителски операции:

\begin{itemize}
	\item \textit{Потребителят иска да зареди мрежа, съхранена в базата данни, и да стартира симулация}
	
	\begin{enumerate}
		\item Потребителят отваря клиентското приложение.
		\item Клиентското приложение изпраща заявка за определена мрежа към сървъра.
		\item Сървърът изпраща съответната заявка към базата данни.
		\item Базата от данни получава заявката и връща съответния запис.
		\item Сървърът препраща получения запис към клиентското приложение.
		\item Клиентът визуализира мрежата.
		\item Клиентското приложение стартира логиката за визуализация на стъпките на симулацията.
	\end{enumerate}
	
	
	\item\textit{Потребителят иска да съхрани мрежа в базата данни}
	\begin{enumerate}
		\item Клиентското приложение изпраща заявка към сървъра с мрежата, която иска да съхрани.
		\item Сървърът изпраща съответната заявка към базата данни за създаване на нов запис.
		\item При успех сървърът връща на клиента адрес, чрез който може да достъпи новия запис.
		\item Клиентът визуализира адреса на потребителя.
	\end{enumerate}
	
	\item \textit{Сървърът периодично изтрива неизползваните мрежи}
	\begin{enumerate}
		\item Сървърът предварително задава период, през който да се изпълнява действието.
		\item Когато настъпи време за изпълнение, сървърът заявява всички записи, които отговарят на критериите за неизползваемост.
		\item В случай че има такива, сървърът изисква да бъдат изтрити от базата.
	\end{enumerate}
	
 \end{itemize}


\subsubsection{Физическа/\textit{Deployment} архитектура и облачна инфраструктура}  

\hspace*{2.5em}Внедряването на приложението се реализира чрез два облачни доставчика -- \textbf{\textit{Render}} (\url{https://render.com/}) и \textbf{\textit{Atlas}} (\url{https://cloud.mongodb.com/}).

\textbf{\textit{Render}} е облачна платформа за хостинг и внедряване на уеб приложения, сървъри, статични сайтове и контейнери. Тя предлага автоматизирани процеси за внедряване, работещи директно с \textit{Git} хранилищата на разработчиците, което улеснява \textit{CI/CD} процесите. \textbf{\textit{Render}} предоставя разнообразни услуги, включително уеб хостинг, статични сайтове, бекенд услуги и управление на контейнери. Отличава се с автоматично скалиране и лесен за използване интерфейс, което позволява бързо внедряване и управление на приложения \cite{RenderWebServices}.

\textbf{\textit{MongoDB Atlas}} е управлявана облачна услуга за бази данни, предоставяна от \textit{MongoDB}. Тя предлага напълно автоматизирано управление на инфраструктурата и основните операции по базата данни, като осигуряване, конфигуриране, мащабиране и мониторинг. \textbf{\textit{Atlas}} включва вградени механизми за висока достъпност, резервни копия и сигурност, което позволява на разработчиците да се фокусират върху създаването на функционалности, вместо върху административни задачи \cite{MongoDBAtlasDocs}.

Компонентите на системата са разположени на различни места в облачната среда (\textbf{Фигура~\ref{fig:cloud-deployment}}). Инстанцията на сървъра се поддържа в инфраструктурата на облачния доставчик \textbf{\textit{Render}}. Също такa, сървърът е отговорен за доставката на клиентското приложение до потребителя при поискване от браузъра на клиента. Зареденото клиентско приложение комуникира със сървъра чрез сигурния протокол \textit{HTTPS}. Инстанцията на базата данни се поддържа в инфраструктурата на облачния доставчик \textbf{\textit{Atlas}}. Сървърът комуникира с базата данни по сигурния протокол \textit{HTTPS}, като при създаване на връзка към базата се удостоверява с потребителско име и парола.


\vspace{0.25cm}

\begin{figure}[H]
	\centering
    \captionsetup{width=0.8\textwidth}
	\includegraphics[width=0.8\linewidth]{"3. onlineGn/images/cloud-deployment.png"}
	\caption{Диаграма, показваща топологичното разположение и комуникацията между компонентите на системата (клиентски браузър, \textit{Render} сървър и \textit{MongoDB Atlas} база данни) в облачната среда.}
	\label{fig:cloud-deployment}
\end{figure}


\subsubsection{Технологичен стек и външни зависимости}  

\hspace*{2.5em}За разработването на трите компонента на системата (база данни, сървър, клиент) са използвани различни технологии.

\paragraph{База данни.}
\mbox{}\\
Данните се съхраняват в документи в \textit{MongoDB} (версия 8.0). \textit{MongoDB} е документно--ориентирана база данни с динамична схема. Документите следват \textit{JSON} структура, но физически се съхраняват в бинарния формат \textit{BSON}.

\paragraph{Сървър.}
\mbox{}\\
За разработката на сървърната част е използван езикът \textit{JavaScript}, работещ в средата за изпълнение \textit{Node.js}. За реализиране на \textit{REST API} слоя е използвана библиотеката \textit{Express.js}, тъй като предоставя стандартизиран механизъм за маршрутизация. Това позволява заявките да се обработват по единен начин чрез централизирано обработване на входа, валидация, логване, унифицирана обработка на грешки и модулна организация на маршрутите, което намалява дублирането на код и риска от непоследователно поведение между отделните крайни точки. 
Проектът е разработен с \textit{Node.js} 22.x (клон \textit{„Jod“}), избран като \textit{LTS} линия поради предвидимия жизнен цикъл и дългосрочните актуализации по сигурност. Версия 22.x преминава в \textit{Active LTS} на 29.10.2024 г. и остава под поддръжка до 04.2027 г., което я прави подходяща за продукционна употреба и намалява риска от принудителни миграции в кратък период. 
В хода на разработката са следвани утвърдени практики (ясна слоеста структура на кода, ограничаване на зависимостите от недокументирани/нестабилни \textit{API}, централизиран \textit{error-handling} и предвидима конфигурация), което поддържа ниска цена на миграция при преминаване към следваща \textit{LTS} версия.\\

Базата данни и сървърът се използват за решението на \ref{task:share-storage}, като осигуряват надеждна инфраструктура за централизирано съхранение на ОМ моделите и генериране на уникални хипервръзки за тяхното споделяне.

По-долу са изброени библиотеките, които са използвани в разработката на приложението (самите библиотеки и техните версии са изброени в \textit{package.json} файла на проекта).

\paragraph{Външни библиотеки, от които зависи сървърната част.}

\begin{itemize}
	\item \texttt{\textbf{"\textit{body-parser}"}} -- предоставя \textit{middleware} (междинен софтуерен слой), който служи за прочитане и обработване на тялото (\textit{body}) на \textit{HTTP} заявката.
	
	\item \texttt{\textbf{"\textit{cors}"}} -- библиотека, която позволява на сървъра да разрешава или забранява заявки от други домейни -- \textit{Cross-Origin Resource Sharing} (\textit{CORS}).
	
	\item \texttt{\textbf{"\textit{cron}"}} -- библиотека, която позволява насрочването на задачи в собствен \textit{Node.js} сървър, които да бъдат изпълнявани на база разписание/интервал.
	
	\item \texttt{\textbf{"\textit{express}"}} -- библиотека, която улеснява писането на уеб сървъри и \textit{API}-та с \textit{Node.js}.
	
	\item \texttt{\textbf{"\textit{is-svg}"}} -- библиотека за проверка дали даден текст е валиден \textit{SVG} код.
	
	\item \texttt{\textbf{"\textit{mongoose}"}} -- \textit{ODM} (\textit{Object Document Mapper}), предоставящ високо ниво на абстракция при работа с \textit{MongoDB}. Позволява дефиниране на модели и схеми, валидиране на данни и улеснява управлението на колекциите.

\end{itemize}

\paragraph{Библиотеки, нужни само за процеса на разработване.}
~\\
\vspace{-0.85cm}

\begin{itemize}
	\item \texttt{\textbf{"\textit{nodemon}"}} -- инструмент за \textit{Node.js}, който автоматично рестартира приложението, когато засече промени във файловете на проекта по време на разработка.
\end{itemize}

\vspace{-0.75cm}

\paragraph{Клиент.}
\mbox{}\\
За разработването на приложението e използвана работната рамка за разработка на уеб приложения \textbf{\textit{Angular}} \cite{AngularDocs}. \textit{Angular} е създаден и поддържан от \textit{Google} и се използва за изграждане на едностранични приложения (\textit{Single Page Applications -- SPA}), които работят без презареждане на страницата. Рамката предлага стабилна екосистема за разработка на динамични и гъвкави уеб приложения, включително двупосочно свързване на данни, инжектиране на зависимости и модулна архитектура. Освен това много основни функционалности са вградени и се поддържат директно от екипа на \textit{Angular}, което намалява нуждата от външни библиотеки.

\textit{Angular} е изграден върху \textit{TypeScript} -- силно типизирано разширение на \textit{JavaScript}, което подобрява качеството на кода, поддръжката и продуктивността на разработчиците чрез статично типизиране \cite{TypeScriptHandbook}. При разработката на приложението е използван стиловият формат \textit{SCSS}, който предоставя разширени възможности като променливи, вложени правила и \textit{mixins}.

Използваната версия на \textit{Angular} е 16, тъй като към началото на разработката на приложението това беше текущата версия с дългосрочна поддръжка (\textit{LTS}). По време на целия процес са спазени утвърдените добри практики за разработка, което гарантира лесно и безпроблемно преминаване към следващата \textit{LTS} версия в бъдеще.


\paragraph{Библиотеки, от които зависи клиентското приложение (самите библиотеки и техните версии са изброени в \textit{package.json} файла на проекта).}
\mbox{} \\

\vspace{-0.4cm}

За начало, клиентската част зависи от всички пакети, които идват с рамката \textit{Angular}:

\vspace{-0.4cm}

\begin{itemize}
	\item \texttt{\textbf{"\textit{@angular/animations}"}} -- Осигурява поддръжка за добавяне на анимации в \textit{Angular} приложения.
	\item \texttt{\textbf{"\textit{@angular/common}"}} -- Съдържа често използвани функции на \textit{Angular} (като пайпове, директиви и услуги).
	\item \texttt{\textbf{"\textit{@angular/compiler}"}} -- Позволява компилация на \textit{Angular} шаблони и компоненти.
	\item \texttt{\textbf{"\textit{@angular/core}"}} -- Основни функционалности за \textit{Angular} приложения, включително компоненти, основни модули и помощни функции.
	\item \texttt{\textbf{"\textit{@angular/forms}"}} -- Поддържа изграждане на форми и управление на контроли на форми в \textit{Angular} приложения.
	\item \texttt{\textbf{"\textit{@angular/platform-browser}"}} -- Включва необходими услуги и инструменти за изпълнение на \textit{Angular} приложения в браузъра.
	\item \texttt{\textbf{"\textit{@angular/platform-browser-dynamic}"}} -- Поддържа компилацията на \textit{Angular} приложения в браузъра по време на изпълнение.
	\item \texttt{\textbf{"\textit{@angular/router}"}} -- Управлява маршрутизацията и навигацията в \textit{Angular} приложения. 
	\item \texttt{\textbf{"\textit{rxjs}"}} -- Библиотека за реактивно програмиране, използваща наблюдаеми обекти (\textit{observables}) и позволяваща асинхронно програмиране.
	\item \texttt{\textbf{"\textit{tslib}"}} -- Библиотека за изпълнителната среда, която разполага с всички помощни \textit{TypeScript} функции.
	\item \texttt{\textbf{"\textit{zone.js}"}} -- Управлява контекста на изпълнение за асинхронни операции, използва се за откриване на промени в данните от \textit{Angular}.
\end{itemize}


\textbf{Допълнително инсталирани библиотеки:}

\vspace{-0.3cm}

\begin{itemize}
	\item \texttt{\textbf{"\textit{@angular/cdk}"}} -- Комплект за разработка на компоненти (\textit{Component Development Kit}) предоставя общи градивни елементи за създаване на компоненти в \textit{Angular}.
	\item \texttt{\textbf{"\textit{@angular/material}"}} -- \textit{Angular} компоненти, базирани на \textit{Material Design}, предоставящи \textit{UI} елементи като бутони, карти и други.
	\item \texttt{\textbf{"\textit{ngx-toastr}"}} -- Библиотека за показване на персонализируеми toast известия в \textit{Angular} приложения, предоставяща ненатрапчиви сигнали или съобщения към потребителите.
	\item \texttt{\textbf{"\textit{d3}"}} -- Мощна библиотека за създаване на визуализации, базирани на данни, която позволява динамична и интерактивна графика чрез \textit{HTML}, \textit{SVG} и \textit{CSS} \cite{D3Transition}. Тази библиотека е нужна за визуализация на цялата ОМ, анимиране на симулациите и внедряване на персонализирани взаимодействия с нея.
	\item \texttt{\textbf{"\textit{d3-graphviz}"}} -- Допълнително разширение за \textit{D3.js}, което предоставя инструменти за рендиране на диаграми и графи на езика \textit{Graphviz DOT} директно в браузъра. Този пакет беше необходим, защото алгоритъмът за персонализирана визуализация генерира \textit{DOT} низ.
\end{itemize}

\vspace{0.35cm}

\subsection{Дизайн и модули на приложението}
\hspace*{2.5em}В контекста на \ref{task:architecture}, тук се представя детайлната модулна организация на приложението \textit{OnlineGN}. Описват се основните функционални компоненти (съхранение, визуализация, редакция, симулация, импорт/експорт и \textit{API}), техните специфични отговорности и начините им на взаимодействие. Избраният модулен дизайн цели постигане на висока степен на разделимост на отговорностите (\textit{separation of concerns}), улеснена поддръжка и мащабируемост на системата.

\subsubsection{Модул за устойчивост на данните (\textit{GenNet Persistence})}
\label{subsec:module-persistence}

\hspace*{2.5em}\textit{GenNet Persistence} реализира слоя за устойчивост на данните в системата и осигурява надеждно съхранение и последващо извличане на моделите на ОМ. Основната му роля е да отдели тази логика от останалите функционални модули, като предостави единен механизъм за запис, зареждане и управление на жизнения цикъл на записите (например автоматично изчистване на неизползвани мрежи). Декомпозицията на модула е показана на \textbf{Фигура~\ref{fig:module-persistence-decomposition}}.
\begin{figure}[H]
	\centering
	\includegraphics[width=0.5\linewidth]{"3. onlineGn/images/module-persistence-decomposition.png"}
	\captionsetup{width=0.7\textwidth}
	\caption{Декомпозиция на модула за устойчивост на данните (\textit{GenNet Persistence}).}
	\label{fig:module-persistence-decomposition}
\end{figure}

\begin{itemize}
	\item Подмодулът за зареждане на мрежа (\textit{Loading}) се грижи за функционалността да зареди мрежа по даден уникален идентификатор. При всяко зареждане се актуализира и полето кога последно записът е бил достъпен.
	
	\item Подмодулът за записване на мрежа (\textit{Saving}) се грижи за запазване на данните за една мрежа, така че да може да бъде прочетена на по-късен етап.
	
	\item Подмодулът за изтриване на изтекли мрежи (\textit{Expiration}) се грижи за изтриването на мрежите, които не са били достъпвани в последните 90 дни.
\end{itemize}

\subsubsection{\textit{JSON} формат за описание на обобщени мрежи}
\label{subsec:gn-json-format}

\hspace*{2.5em}Разработеният \textit{JSON} формат предоставя унифицирано, машинночетимо представяне на ОМ, пригодено за директна интерпретация от симулатора. Форматът на данните, използвани за записване и прочитане на информацията за една мрежа, е дефиниран от интерфейса \textit{GenNetExportModel}. Един документ описва както \emph{структурата} на мрежата (позиции, преходи и връзки между тях), така и \emph{началната конфигурация} на ядрата (tokens), параметри на изпълнение, набор от потребителски дефинирани функции, използвани при симулацията, както и предварително генерирана визуализация на мрежата. Форматът е организиран в четири основни секции: \texttt{genNetRaw}, \texttt{settings}, \texttt{code} и \texttt{svg}.

\paragraph{Коренова структура.}\mbox{}\\
Кореновият обект съдържа:
\begin{itemize}
\item \texttt{\textbf{\textit{genNetRaw}}} -- обект от тип \textit{GenNetRaw}, съдържащ декларативно описание на ОМ;
\item \texttt{\textbf{\textit{settings}}} -- обект от тип \textit{GenNetSettings}, съдържащ настройки на симулационната среда;
\item \texttt{\textbf{\textit{code}}} -- низ, съдържащ написан от потребителя \textit{JavaScript} код, който да се компилира от клиентското приложение, за да може да бъде използван в необходимите ситуации;
\item \texttt{\textbf{\textit{svg}}} -- низ, който съдържа предварително генерирана визуализация на мрежата. Възможно е стоността да е празна.
\end{itemize}
Така се постига ясно разграничение между \emph{структурни данни} и \emph{поведенческа спецификация}.

\paragraph{Секция \texttt{genNetRaw}.}\mbox{}\\
Секцията \texttt{genNetRaw} представлява основните данни за мрежата и съдържа три основни колекции: \texttt{places}, \texttt{transitions} и \texttt{tokens}.

\begin{itemize}
\item \texttt{\textbf{\textit{places}}} -- масив от обекти от тип \textit{Place}.
\begin{itemize}
\item \texttt{\textbf{\textit{id}}} -- низ, представляващ уникален идентификатор на позицията;
\item \texttt{\textbf{\textit{name}}} -- низ, представляващ име на позицията;
\item \texttt{\textbf{\textit{charFunc}}} -- низ или празна стойност (\texttt{null}), представляващ име на функцията, която да се използва за трансформиране на характеристиките на ядрото;
\item \texttt{\textbf{\textit{mergeFunction}}} -- низ или празна стойност (\texttt{null}), представляващ име на функцията, която да се използва за сливане на ядра. Функцията приема масив от ядра и връща масив от ядра, за да позволи сливане на $n$ ядра в $k$ ядра ($k \leq n$);
\item \texttt{\textbf{\textit{priority}}} -- число, представляващо приоритета на позицията;
\item \texttt{\textbf{\textit{capacity}}} -- число, представляващо колко ядра могат да бъдат едновременно на позицията.
\end{itemize}
\end{itemize}

Пример за обект от тип \textit{Place} е показан във \textbf{\autoref{lst:place-example}}.
\begin{lstlisting}[
	language=JavaScript,
	caption={Пример за \textit{JSON} представяне на обект от тип \textit{Place}.},
	label={lst:place-example}
	]
{
    "id": "p4",
    "name": "p_4",
    "charFunc": null,
    "mergeFunction": "mergeTokensOnP4",
    "priority": 2,
    "capacity": 1
}
\end{lstlisting}

\vspace{-10pt}
\begin{itemize}
\item \texttt{\textbf{\textit{transitions}}} -- масив от обекти от тип \textit{Transition}. Ако преходът има $n$ позиции, които влизат в него, и $m$ позиции, които излизат от него, то масивът \texttt{data} трябва да има $n \cdot m$ записа.
\begin{itemize}
\item \texttt{\textbf{\textit{id}}} -- низ, представляващ уникален идентификатор на прехода;
\item \texttt{\textbf{\textit{name}}} -- низ, представляващ име на прехода;
\item \texttt{\textbf{\textit{splitFunction}}} -- низ или празна стойност (\texttt{null}), представляващ име на функцията, която да се използва за определяне на това дали да се раздели ядро;
\item \texttt{\textbf{\textit{priority}}} -- число, представляващо приоритета на прехода;
\item \texttt{\textbf{\textit{data}}} -- масив от обекти от тип \textit{TransitionItem}, съдържащ информация за връзките между позициите:
\begin{itemize}
\item \texttt{\textbf{\textit{source}}} -- низ, представляващ \texttt{\textbf{\textit{id}}} на начална позиция;
\item \texttt{\textbf{\textit{target}}} -- низ, представляващ \texttt{\textbf{\textit{id}}} на крайна позиция;
\item \texttt{\textbf{\textit{predicate}}} -- низ или празна стойност (\texttt{null}), представляващ име на функцията, която да се използва за определяне дали ядро може да премине по тази връзка;
\item \texttt{\textbf{\textit{capacity}}} -- число, представляващо колко ядра могат да минат едновременно по тази връзка.
\end{itemize}
\end{itemize}
\end{itemize}

Пример за обект от тип \textit{Transition} е показан във \textbf{\autoref{lst:transition-example}}.
\begin{lstlisting}[
	language=JavaScript,
	caption={Пример за \textit{JSON} представяне на обект от тип \textit{Transition}.},
	label={lst:transition-example}
	]
{
    "id": "T1",
    "name": "T_1",
    "splitFunction": "splitTokenOnT1",
    "priority": 2,
    "data": [
        {
            "source": "p1",
            "target": "p3",
            "predicate": null,
            "capacity": 1
        }
    ]
}
\end{lstlisting}

\vspace{-10pt}
\begin{itemize}
\item \texttt{\textbf{\textit{tokens}}} -- масив от обекти от тип \textit{Token}, дефиниращ началната конфигурация на симулацията.
\begin{itemize}
\item \texttt{\textbf{\textit{id}}} -- низ, представляващ уникален идентификатор на ядрото;
\item \texttt{\textbf{\textit{name}}} -- низ, представляващ име на ядрото;
\item \texttt{\textbf{\textit{host}}} -- низ, представляващ \texttt{\textbf{\textit{id}}} на текущата позиция на ядрото;
\item \texttt{\textbf{\textit{priority}}} -- число, представляващо приоритет на ядрото;
\item \texttt{\textbf{\textit{initialChars}}} -- обект, съдържащ първоначалните характеристики на ядрото.
\end{itemize}
\end{itemize}

Пример за обект от тип \textit{Token} е показан във \textbf{\autoref{lst:token-example}}.
\begin{lstlisting}[
	language=JavaScript,
	caption={Пример за \textit{JSON} представяне на обект от тип \textit{Token}.},
	label={lst:token-example}
	]
{
    "id": "token1",
    "name": "token_1",
    "host": "p1",
    "priority": 1,
    "initialChars": {
        "a": 1,
        "b": {
            "prop": 1
        },
        "c": "2025-04-12T09:37:16.950Z",
        "d": [
            1, 2
        ]
    }
}
\end{lstlisting}

\vspace{-10pt}
\paragraph{Секция \textit{\texttt{settings}}.}\mbox{}\\
Секцията \textit{\texttt{settings}} съдържа параметри, влияещи на изпълнението, без да променят структурния модел.

\begin{itemize}
\item \texttt{\textbf{\textit{simulationSettings}}} -- обект от тип \textit{SimulationSettings}.
\begin{itemize}
\item \texttt{\textbf{\textit{stepDurationMS}}} -- число, което определя продължителността на една симулационна стъпка в милисекунди.
\end{itemize}
\end{itemize}

Пример за обект от тип \textit{GenNetSettings} е показан във \textbf{\autoref{lst:gennetsettings-example}}.
\begin{lstlisting}[
	language=JavaScript,
	caption={Пример за \textit{JSON} представяне на обект от тип \textit{GenNetSettings}.},
	label={lst:gennetsettings-example}
	]
{
    "simulationSettings": {
        "stepDurationMS": 5000
    }
}
\end{lstlisting}

\vspace{-10pt}
\paragraph{Секция \textit{\texttt{code}}.}\mbox{}\\
Секцията \textit{\texttt{code}} съдържа низ с потребителски \textit{JavaScript} код, който дефинира функции и променливи, използвани при симулацията. За да бъде кодът достъпен в контекста на симулатора, на най-външно ниво трябва да бъде дефиниран обектът \texttt{\textbf{\textit{this.globalObj}}}, който служи като пространство от имена за потребителските дефиниции. Ако той не бъде дефиниран или някое от задължителните му полета не бъде дефинирано, те се добавят по подразбиране от симулатора. В \texttt{\textbf{\textit{this.globalObj}}} се съдържат посочените по-долу елементи.

\begin{itemize}
\item \texttt{\textbf{\textit{tokens}}} -- задължително поле в \texttt{\textbf{\textit{this.globalObj}}}, което трябва да реферира \texttt{\textbf{\textit{this.tokenMover.tokens}}} и осигурява достъп до текущата колекция от ядра;
\item \texttt{\textbf{\textit{functions}}} -- функции, дефинирани в \texttt{\textbf{\textit{this.globalObj}}}, които могат да бъдат извиквани чрез имената си от съответните полета в \texttt{\textbf{\textit{genNetRaw}}}.
\item \texttt{\textbf{\textit{globals}}} -- допълнителни стойности, дефинирани в \texttt{\textbf{\textit{this.globalObj}}}, които могат да се използват като спомагателен контекст при изпълнение на симулацията.
\end{itemize}

Пример за правилно форматиран низ, е показан във \textbf{\autoref{lst:global-obj-example}}.
\begin{lstlisting}[
		language=JavaScript,
		caption={Пример за дефиниране на \textit{globalObj}.},
		label={lst:global-obj-example}
		]
this.globalObj = {
    // reference to all tokens/cores - mandatory
    tokens: this.tokenMover.tokens,
    
    // example of a function that gets a token/core by name
    getTokenByName: function (name) {
        return this.tokens.find((token) => token.name === name);
    },

    // here the user can now define
    // their own variables and functions
};
\end{lstlisting}

Така форматът позволява декларативно задаване на мрежа, като конкретното поведение и условия се дефинират централизирано и се свързват чрез имената на функциите (\textbf{\autoref{lst:gn-json-minimal}}).

\begin{lstlisting}[language=JavaScript,caption={Примерен \textit{JSON} модел на обобщена мрежа с две позиции, един преход и едно ядро.},label={lst:gn-json-minimal}]
{
  "genNetRaw": {
    "places": [
      {
        "id": "p1",
        "name": "start",
        "charFunc": null,
        "mergeFunction": null,
        "priority": 1,
        "capacity": 1
      },
      {
        "id": "p2",
        "name": "end",
        "charFunc": null,
        "mergeFunction": null,
        "priority": 1,
        "capacity": 1
      }
    ],
    "transitions": [
      {
        "id": "T1",
        "name": "move",
        "splitFunction": null,
        "data": [
          {
            "source": "p1",
            "target": "p2",
            "predicate": "alwaysTrue",
            "capacity": 1
          }
        ],
        "priority": 1
      }
    ],
    "tokens": [
      {
        "id": "k1",
        "name": "kernel1",
        "host": "p1",
        "priority": 1,
        "initialChars": {
          "x": 0
        }
      }
    ]
  },
  "settings": {
    "simulationSettings": {
      "stepDurationMS": 300
    }
  },
  "code": "this.globalObj = {\\n  tokens: this.tokenMover.tokens,\\n\\n  alwaysTrue: function () {\\n    return true;\\n  }\\n};\\n"
}
\end{lstlisting}

\vspace{-10pt}
\paragraph{Секция \textit{\texttt{svg}}.}\mbox{}\\
Секцията съдържа \textit{\texttt{SVG}} низ за визуализация на мрежата. Ако такъв не е предоставен, програмата генерира \textit{SVG} визуализация чрез алгоритъма за автоматично изчертаване, описан в \textbf{\ref{sec:auto-drawing}}

%here i put new page command
\newpage
\subsubsection{Модул за графична визуализация (\textit{SVG/Graphviz)}}  

\hspace*{2.5em}Модулът за графично визуализиране включва в себе си всички функционалности, свързани с визуализиране и анимиране на обекти, както и зареждането на графичната визуализация от различни обекти (\textbf{Фигура~\ref{fig:module-visualization-decomposition}}).

\begin{figure}[H]
	\centering
	\includegraphics[width=0.7\linewidth]{"3. onlineGn/images/module-visualization-decomposition.png"}
	\captionsetup{width=0.8\textwidth}
	\caption{Декомпозиция на модула за графична визуализация (Visualization).}
	\label{fig:module-visualization-decomposition}
\end{figure}

\begin{itemize}
	\item Подмодулът за зареждане на графика от (\textit{Load from}) с неговите подмодули се грижи за визуализация на интерактивната мрежа. В случай че е предоставено предварително \textit{SVG} файл на мрежата, той бива директно зареден от \textit{SVG} модулът, в противен случай бива предоставен \textit{Graphviz} низ, за който подмодулът се грижи да го визуализира като \textit{SVG}.
	
	\item Подмодулът за рисуване (\textit{Draw}) се грижи за рисуването и премахването на различни фигури в мрежата, както и прерисуване на елементи в случай на промяна на техните координати.
	
	\item Подмодулът за анимации (\textit{Animation}) се грижи за всички функционалности, които са свързани с анимиране на графичните елементи. Включва анимирана промяна на позиция и трансформация.
	
	\item Подмодулът за обработване на събития (\textit{Event Handling}) се грижи в случай на взаимодействие на потребителя с графичната визуализация мрежата (например клик с мишката или преминаване с мишката върху елемент) това да се регистрира като събитие и да се уведомят слушателите на това събитие за неговото възникване, за да може да се изпълнят съответните функционалности регистрирани за това събитие.
	
	\item Подмодулът за запазване в \textit{SVG} формат (\textit{Save to SVG}) се грижи за съхранението на текущото състояние на визуализираната мрежа (с нейните визуални промени) в \textit{SVG} формат, така че да може при последващо зареждане на \textit{SVG} файла, мрежата да изглежда по същия начин, както когато е била съхранена. 
	
\end{itemize}

\subsubsection{Модул за симулация на движението на ядрата}

\hspace*{2.5em}Модулът за симулация на движението на ядрата (\textit{Simulation}) се грижи за логиката за правилното изпълнение на симулацията (\textbf{Фигура~\ref{fig:module-simulation-decomposition}}). Също така комуникира с модула за визуализацията, давайки инструкции за преместването на ядрата.

\begin{figure}[H]
	\centering
	\includegraphics[width=0.6\linewidth]{"3. onlineGn/images/module-simulation-decomposition.png"}
	\captionsetup{width=0.7\textwidth}
	\caption{Декомпозиция на модула за симулация на движението на ядрата (Simulation).}
	\label{fig:module-simulation-decomposition}
\end{figure}

\begin{itemize}
	\item Подмодулът за актуализация на визуализацията (\textit{Update visualization}) се използва от останалите два подмодула за извикване на подмодула за анимация (\textit{Animation}) на модула за визуализация (\textit{Visualization}). Неговата задача е да подготви и предаде необходимите данни, за да се покаже коректната анимация при преместване на ядра в потребителския интерфейс.
	
	\item 	Подмодулът за симулация на единичен ход (\textit{Single move}) се грижи за изпълнението на един ход от симулацията. Това включва в себе си правилната обработка и движение на всяко от ядрата за този ход на база правилата на ОМ. След като логически обработи данните, извиква другия подмодул за актуализация на визуализацията (\textit{Update visualization}), който да се погрижи за отразяването на промените и във визуализацията.
	
	\item Подмодулът за пълна симулация (\textit{Full simulation}) има за цел да изпълни всички възможни ходове на симулацията, стартирайки от текущото състояние на мрежата, докато има възможни ходове. За целта той използва подмодулa за симулация на единичен ход (\textit{Single move}) за изпълнението на всеки ход.
	
	\item Подмодулът за рестартиране на симулацията (\textit{Reset}) осигурява възможност за връщане на симулацията към първоначалното ѝ състояние. Той комуникира с подмодула за актуализация на визуализацията (\textit{Update visualization}), за да се синхронизира визуалното представяне със състоянието на симулацията.
	
\end{itemize}

%here i put new page command
\newpage
\subsubsection{\textit{API} и скриптов интерфейс}\label{sec:api-script-interface}
%\mbox{}\\
Сървърът предоставя две крайни точки за достъп от клиент:

\vspace{-10pt}
\begin{itemize}
	\item \texttt{\textbf{\textit{GET /api/genNets/:id}}}, където \texttt{\textbf{\textit{id}}} е идентификатор на мрежата
	\item връща \texttt{\textbf{\textit{status code 200}}} и мрежата с този идентификатор, ако тя бъде намерена
	\item връща \texttt{\textbf{\textit{status code 404}}}, ако мрежата не бъде намерена
	\item в случай, че мрежата е намерена се връща обект със следните полета:
	\begin{itemize}
		\item \texttt{\textbf{\textit{id}}} -- идентификатор на мрежата
		\item \texttt{\textbf{\textit{content}}} -- обект от тип \textit{GenNetExportModel}, описан по-горе
		\item \texttt{\textbf{\textit{lastFetchedAt}}} -- дата и час кога за последно е била достъпвана мрежата
	\end{itemize}
\end{itemize}

Пример за \textit{GET} заявка към \textit{API} за \textit{GenNet} е показан във \textbf{\autoref{lst:api-get-gennet}}.
\begin{lstlisting}[
	language=BASH,
	caption={Пример за \textit{GET} заявка към \textit{API} за \textit{GenNet}.},
	label={lst:api-get-gennet}
	]
GET /api/genNets/67fd54e71d8d07452f305ecc
\end{lstlisting}

Примерен \textit{JSON} отговор от \textit{API} за \textit{GenNet} е показан във \textbf{\autoref{lst:api-response-gennet}}.
\begin{lstlisting}[
	language=JavaScript,
	caption={Примерен \textit{JSON} отговор от \textit{API} за \textit{GenNet}.},
	label={lst:api-response-gennet}
	]
{
    "content": {...},
    "id": "67fd54e71d8d07452f305ecc",
    "lastFetchedAt": "2025-08-29T07:57:00.632Z"
}
\end{lstlisting}

\vspace{-10pt}
\begin{itemize}
	\item \texttt{\textbf{\textit{POST /api/genNets}}} -- записва подадена мрежа в базата от данни
	\item В тялото на заявката трябва да се постави обект с едно поле \textit{content} представляващо обект от тип \textit{GenNetExportModel}
	\item При успешно записване се връща \texttt{\textbf{\textit{status code 201}}} и също така се връща обект с едно поле \texttt{\textbf{\textit{id}}}, представляващо уникалния идентификатор на мрежата. \texttt{\textbf{\textit{id}}} се генерира от базата данни при записване.
	\item При неуспешно записване се връща \texttt{\textbf{\textit{status code 400}}} и съответната грешка
\end{itemize}

Примерен \textit{JSON} за изпращане на заявка е показан във \textbf{\autoref{lst:api-request-example}}.
\begin{lstlisting}[
	language=JavaScript,
	caption={Примерен \textit{JSON} за изпращане на заявка.},
	label={lst:api-request-example}
	]
	{
		content: {....}
	}
\end{lstlisting}

Примерен \textit{JSON} отговор от \textit{API} е показан във \textbf{\autoref{lst:api-response-example}}.
\begin{lstlisting}[
	language=JavaScript,
	caption={Примерен \textit{JSON} отговор от \textit{API}.},
	label={lst:api-response-example}
	]
	{
		id: "68b15de499090443949c3aa7"
	}
\end{lstlisting}

\paragraph{Изпълнение на характеристични функции и предикати.}
\mbox{}\\
В рамките на скриптовия интерфейс характеристичните функции и предикатите могат да бъдат задавани като \textit{JavaScript} код, който се съхранява като част от описанието на мрежата и се изпълнява в клиентската част на приложението. Поради това, ако в такава функция или предикат бъде въведен недоброжелателен код, той би могъл да повлияе върху правилното изпълнение на клиентската програма или на конкретната заредена мрежа. В същото време в приложението не се съхраняват данни за идентификация на потребител нито на сървъра, нито в клиентската среда, а данните в клиента са свързани единствено с текущо заредената мрежа. Поради тази причина евентуално недоброжелателно съдържание в характеристични функции или предикати не води до компрометиране на потребителски данни. Освен това низът, съдържащ съответния код, се записва като стойност и не модифицира заявката за запис към базата от данни.

\vspace{0.2cm}

\subsubsection{Модул за импортиране и експортиране на мрежи}

\hspace*{2.5em}Модулът за импортиране и експортиране на мрежи се грижи за зареждането и запазването в различни формати на данните на мрежата. Форматът на данните е описан в \textbf{точка~\ref{sec:api-script-interface}} 

Симулаторът предоставя възможност за запазване на мрежа:

\vspace{-0.4cm}

\begin{itemize}
	\item в \textit{JSON} файл
	\item в \textit{TeX} формат
	\item в \textit{JSON} формат в база от данни
\end{itemize}

Съответно, симулаторът може да зареди мрежа:

\vspace{-0.4cm}

\begin{itemize}
	\item от файл
	\item от съхранените данни в \textit{JSON} формат в базата от данни
\end{itemize} 

Декомпозицията на модула за импортиране и експортиране на данните на мрежата е показана на \textbf{Фигура~\ref{fig:data-import-export-decomposition}}.

\begin{figure}[H]
	\centering
	\includegraphics[width=0.7\linewidth]{"3. onlineGn/images/data-import-export-decomposition.png"}
	\captionsetup{width=0.81\textwidth}
	\caption{Декомпозиция на модула за импортиране и експортиране на данните на мрежата (\textit{Data Import/Export}).}
	\label{fig:data-import-export-decomposition}
\end{figure}


Модулът за импортиране и експортиране на данните на мрежата (\textit{Data Import/Export}) се разделя на два подмодула: \textit{импортиране (Import)}, грижещ се за зареждане на мрежа от различни формати, и \textit{експортиране (Export)}, грижещ се за да запише данните на мрежата в различни формати.

\paragraph{Подмодул \textit{Import}.}
\mbox{}\\
Подмодулът \textit{Import} се разделя на два подмодула:
\begin{itemize}
	\item \textbf{Подмодул \textit{JSON}} -- Отговаря за прочитане данните на мрежа от \textit{JSON} формат, както и валидиране на данни, предоставени в \textit{JSON} формат, спрямо интерфейса, дефиниран в \textbf{точка~\ref{sec:api-script-interface}} При успешна валидация се създават необходимите обекти за визуализация на мрежата. Този подмодул се използва също и за зареждане на данни от сървър, тъй като те пристигат на клиента във формат \textit{JSON}.
	
	\item \textbf{Подмодул \textit{File}} -- Отговаря за обработката на файлове, предоставени от потребителя. Ако файлът е валиден \textit{JSON}, съдържанието се предава на \textit{JSON} подмодула за по-нататъшна обработка. 
	
\end{itemize}

\paragraph{Подмодул \textit{Export}.}
\mbox{}\\
Подмодулът \textit{Export} се разделя на три подмодула:
\begin{itemize}
	\item \textbf{Подмодул \textit{TeX}} -- Отговаря за записването на мрежата в \textit{TeX} формат. Приема текущата визуализация на мрежата и я преобразува в \textit{TeX} формат, което позволява интеграцията ѝ в различни научни документи.
	
	\item \textbf{Подмодул \textit{JSON}} -- Отговаря за преобразуване на текущо заредената мрежа в JSON формат, включително данни за ядра и визуализация, спрямо описания \textit{JSON} формат в \textbf{точка~\ref{sec:api-script-interface}} Също така, този подмодул се използва и при съхранение на мрежата в базата от данни, тъй като там тя бива запазена в \textit{JSON} формат.
	
	\item \textbf{Подмодул \textit{File}} -- Отговаря за създаването на файл с информация за мрежата, който може да бъде свален от потребителя. За целта подмодулът \textit{JSON} генерира необходимите данни, които след това се записват във файла.
\end{itemize}


\subsubsection{Модул за визуализация на динамични процеси и анимация}

\hspace*{2.5em}За осигуряване на плавна визуализация на движението на ядрата в рамките на симулатора е използвана библиотеката \textit{d3}. Визуалното представяне на всяко ядро се реализира чрез създаването на \textit{SVG} група (\textit{g element}), обединяваща графичен символ (кръг) и текстова информация.

Основното предизвикателство при визуализирането на ОМ е преместването на тези групи по сложни траектории, дефинирани чрез \textit{SVG path} елементи, чийто атрибут \textit{d} съдържа координатите на точките за смяна на посоката. Процесът на анимация е реализиран чрез метода на интерполация:

\begin{itemize}
	\item \textbf{Инициализация и затваряне:} Използва се функцията \texttt{translateAlong(path)}, която действа като фабрика за функции. При извикването си тя изчислява общата дължина на пътя $l$ чрез \texttt{getTotalLength()} и я капсулира в затваряне (\textit{closure}), за да бъде достъпна по време на целия цикъл на анимацията.
	
	\item \textbf{Интерполация (\textit{Tweening}):} Чрез метода \texttt{attrTween} на библиотеката \textit{d3} се дефинира преход, който предава параметър $t \in [0, 1]$. Този параметър отразява прогреса на анимацията във всеки времеви откъс.

	\item \textbf{Изчисление на позицията:} Във всеки момент от анимацията се извиква вътрешната функция, която използва \texttt{getPointAtLength(t * l)}. Това позволява намирането на точните $(x, y)$ координати върху кривата, съответстващи на текущия етап от прехода.

	\item \textbf{Трансформация:} На база на получените координати, \textit{CSS} атрибутът \textit{translate} на \textit{SVG} групата се актуализира динамично, „пренасяйки“ ядрото по траекторията.
\end{itemize}

След приключване на анимационния цикъл (\textit{end of transition}), елементът се премахва автоматично от \textit{DOM} дървото. Това гарантира освобождаване на системни ресурси и коректно отразява логическото преминаване на ядрото през съответното ребро от мрежата.

\vspace{0.25cm}
%here i put new page command
\newpage

\subsection{Алгоритми и имплементационни детайли}
\hspace*{2.5em}Описват се основните алгоритми и ключови имплементационни решения, използвани при реализирането на системата \textit{OnlineGN}.

\subsubsection{Разработен алгоритъм 1: автоматично изчертаване на обобщени мрежи}

Алгоритъмът за автоматично изчертаване на ОМ е показан във \textbf{\autoref{lst:algorithm-input-interfaces}}.

\begin{lstlisting}[
	language=TypeScript,
	caption={Алгоритъм за автоматично изчертаване на ОМ.},
	label={lst:algorithm-input-interfaces}
	]
export interface AlgorithmInput {
    places: AlgorithmPlace[];
}

export interface AlgorithmPlace {
    name: string;
    beginTransition?: string;
    endTransition?: string;
}

import {
    AlgorithmInput,
    AlgorithmPlace
} from './algorithm-input.model';

function generateInvisibleNodes(
    ss: string[],
    transitionName: string,
    inputsSize: number,
): void {
    // Creating the invisible nodes
    for (let i = 0; i < inputsSize; i++) {
        ss.push('\n');
        ss.push(
            `invis_node_${transitionName}_${i}[color = black, ` +
            `shape = point, width = 0.01, ` +
            `height = 0.01, ` +
            `class="invis_node ` +
            `invis_node_${transitionName}_${i}"];`,
        );
        ss.push('\n');
    }
}

function generateEdgeBetweenInvisibleNodesAndTransition(
    ss: string[],
    transitionName: string,
    inputsSize: number,
): void {
    for (let i = 0; i < inputsSize; i++) {
        ss.push('\n');
        ss.push(
            `invis_node_${transitionName}_${i} -> ` + 
            `${transitionName}:w [color = black, ` + 
            `class="invis_node_${transitionName}_${i}___` + 
            `${transitionName}"];`, 
        );
        ss.push('\n');
    }
}

function generateOutgoingPlacesFromTransition(
    ss: string[],
    destinations: string[],
): void {
    ss.push('{ rank = same; ');
    for (const outputPlace of destinations) {
        ss.push('\n');
        ss.push(`${outputPlace};`);
        ss.push('\n');
    }
    ss.push('}');
    ss.push('\n');
}

function generateOutgoingEdgesFromTransition(
    ss: string[],
    transitionName: string,
    destinations: string[],
): void {
    for (const outputPlace of destinations) {
        ss.push('\n');
        ss.push(
            `${transitionName} -> ${outputPlace} ` +
            `[color = black, class="${transitionName}___` +
            `${outputPlace}"];`,
        );
        ss.push('\n');
    }
}

function generateTransitionsString(
    ss: string[],
    transitionToOutputs: Map<string, string[]>,
    transitionToInputs: Map<string, string[]>,
): string {
    transitionToOutputs.forEach((outputs, transitionName) => {
        const inputsSize: number =
            transitionToInputs.get(transitionName)?.length || 0;
        
        const transitionWidth: number = 0.01;
        
        ss.push('\n');
        ss.push(`subgraph cluster_${transitionName} {`);
        ss.push('\n');
        ss.push('style=invis');
        ss.push('\n');
        ss.push(`subgraph cluster_${transitionName}_0 {`);
        ss.push('\n');
        
        ss.push(
            `${transitionName}[shape=rect, height=${Math.max(` +
            `outputs.length, inputsSize,` +
            `)}` + 
            `, width=${transitionWidth}, ` + 
            `label = "", ` + 
            `fixedsize=true, ` + 
            `class="transition ${transitionName}"];`, 
        );
        ss.push('\n');
            
        generateInvisibleNodes(ss, transitionName, inputsSize);
        
        generateEdgeBetweenInvisibleNodesAndTransition(
            ss,
            transitionName,
            inputsSize,
        );
        
        ss.push('}'); // END cluster of subgraph - 2
        ss.push('\n');
            
        const outputPlaces = 
            transitionToOutputs.get(transitionName) || [];
            
        generateOutgoingPlacesFromTransition(ss, outputPlaces);
        
        generateOutgoingEdgesFromTransition(
            ss, transitionName, outputPlaces
        );
            
        ss.push('}'); // END cluster of subgraph - 1
        ss.push('\n');
    });

    // Joining the string array into a single string
    return ss.join('');
}

function fillMaps(
    genNet: AlgorithmPlace[],
    transitionToOutputs: Map<string, string[]>,
    transitionToInputs: Map<string, string[]>,
    placeToTransition: Map<string, string>,
): void {
    genNet.forEach((place) => {
        if (place.beginTransition !== undefined) {
            const findItOutputs =
                transitionToOutputs.get(place.beginTransition);
            if (findItOutputs) {
                findItOutputs.push(place.name);
            } else {
                transitionToOutputs.set(
                    place.beginTransition,
                    [place.name]
                );
            }
        }
        
        if (place.endTransition !== undefined) {
            const findItInputs = 
                transitionToInputs.get(place.endTransition);
            if (findItInputs) {
                findItInputs.push(place.name);
            } else {
                transitionToInputs.set(
                    place.endTransition,
                    [place.name]
                );
            }
            
            placeToTransition.set(
                place.name, place.endTransition
            );
        }
    });
}

function generateCircleNode(
    res: string[],
    genNet: AlgorithmPlace[]
): void {
    genNet.forEach((place) => {
        res.push('\n');
        res.push(
            `${place.name} [shape = circle, ` +
            `width=.3, label = "", ` +
            `class="place ${place.name}"];`,
        );
        res.push('\n');
    });
}

function getFirstFreeInvisibleNode(
    transStr: string,
    transitionToFreeInvisNodeIndex: Map<string, number>,
): number {
    const invisNodeIndex: number =
        transitionToFreeInvisNodeIndex.get(transStr) || 0;
    
    transitionToFreeInvisNodeIndex.set(
        transStr, invisNodeIndex + 1
    );
    
    return invisNodeIndex;
}

function generateOutputEdgesFromPlaces(
    res: string[],
    placeToTransition: Map<string, string>,
): void {
    const transitionToFreeInvisNodeIndex: Map<string, number> =
        new Map();
    
    placeToTransition.forEach((transition, place) => {
        const invisNodeIndex: number = getFirstFreeInvisibleNode(
            transition,
            transitionToFreeInvisNodeIndex,
        );
        
        res.push('\n');
        res.push(
            `${place} -> ` +
            `invis_node_${transition}_${invisNodeIndex} ` +
            `[arrowhead=none, color=black, ` +
            `class="${place}___invis_node_${transition} ` +
            `i${invisNodeIndex}"];`,
        );
        res.push('\n');
    });
}

function isValidInput(
    transitionToOutputs: Map<string, string[]>,
    transitionToInputs: Map<string, string[]>,
    placeToTransition: Map<string, string>,
): boolean {
    for (
        const [transitionName, outputs] of transitionToOutputs
    ) {
        if (!transitionToInputs.has(transitionName)) {
            console.log(
                `Transition ${transitionName} ` +
                `has no input places!`,
            );
            return false;
        }
    }
    
    // Check transitionToInputs
    for (const [transitionName, inputs] of transitionToInputs) {
        if (!transitionToOutputs.has(transitionName)) {
            console.log(
                `Transition ${transitionName} ` +
                `has no output places!`,
            );
            return false;
        }
    }
    
    return true;
}

export function generateGraphvizString(
    input: AlgorithmInput 
): string {
    const transitionToOutputs = new Map<string, string[]>();
    const transitionToInputs = new Map<string, string[]>();
    const placeToTransition = new Map<string, string>();
    
    fillMaps(
        input.places,
        transitionToOutputs,
        transitionToInputs,
        placeToTransition,
    );
    
    if (
        !isValidInput(
            transitionToOutputs, 
            transitionToInputs,
            placeToTransition
        )
    ) {
        return '';
    }
    
    const res: string[] = [
        'digraph G { ',
        'rankdir=LR; ',
        'splines=ortho;',
        '\n',
    ];
        
    generateTransitionsString(
        res, transitionToOutputs, transitionToInputs
    );
        
    generateCircleNode(res, input.places);
    generateOutputEdgesFromPlaces(
        res, placeToTransition
    );
        
    res.push('}');
    
    return res.join('');
}
\end{lstlisting}

%\vspace{0.55cm}
\vspace{0.25cm}

\paragraph{Декларации на функции.}
\mbox{}

Пример за декларация на функцията \texttt{generateInvisibleNodes} е показан във \textbf{\autoref{lst:generate-invisible-nodes}}.
\begin{lstlisting}[
	language=TypeScript,
	caption={Декларация на функция \textit{generateInvisibleNodes}.},
	label={lst:generate-invisible-nodes}
	]
function generateInvisibleNodes(
    ss: string[],
    transitionName: string,
    inputsSize: number,
): void
\end{lstlisting}

Функцията създава допълнителни невидими точки, които се поставят пред преходите. За всяка позиция, свързана с даден преход, се генерира по една такава точка. Тези точки се използват при визуализирането на ребрата, които свързват позиции и преходи. Функцията не връща нищо.

Параметрите на функцията са:

\vspace{-0.3cm}

\begin{itemize}
	\item \texttt{\textbf{\textit{ss}}} -- масив от низове, съдържащи резултатния \textit{Graphviz} низ на части, които ще бъдат конкатенирани на по-късен етап.
	\item \texttt{\textbf{\textit{transitionName}}} -- името на прехода, за който ще бъдат добавени невидимите точки.
	\item \texttt{\textbf{\textit{inputSize}}} -- брой необходими невидими точки
\end{itemize}

\vspace{0.55cm}

Пример за декларация на функцията \texttt{generateEdge\-Between\-Invisible\-Nodes\-And\-Transition} е показан във \textbf{\autoref{lst:generate-edge-invisible-nodes-transition}}.
\begin{lstlisting}[
	language=TypeScript,
	caption={Декларация на функция за генериране на ребра към невидими точки.},
	label={lst:generate-edge-invisible-nodes-transition}
	]
function generateEdgeBetweenInvisibleNodesAndTransition(
    ss: string[],
    transitionName: string,
    inputsSize: number,
): void
\end{lstlisting}

Функцията създава ребра между невидимите точки и преходите и ги добавя в резултатния масив. Не връща нищо.

Параметри на функцията:

\vspace{-0.3cm}

\begin{itemize}
	\item \texttt{\textbf{\textit{ss}}} -- масив от низове, съдържащи резултатния \textit{Graphviz} низ на части, които ще бъдат конкатенирани на по-късен етап.
	\item \texttt{\textbf{\textit{transitionName}}} -- името на прехода, за чиито точки трябва да се създадат ребра
	\item \texttt{\textbf{\textit{inputSize}}} -- брой невидими точки.
\end{itemize}

\vspace{0.55cm}

Пример за декларация на функцията \texttt{generateOutgoingPlacesFromTransition} е показан във \textbf{\autoref{lst:generate-outgoing-places-from-transition}}.
\begin{lstlisting}[
	language=TypeScript,
	caption={Декларация на функция \textit{generateOutgoingPlacesFromTransition}.},
	label={lst:generate-outgoing-places-from-transition}
	]
function generateOutgoingPlacesFromTransition(
    ss: string[],
    destinations: string[],
): void
\end{lstlisting}

Функцията създава върхове за всички подадени позиции и ги добавя към резултатния масив. Не връща нищо.

Параметри на функцията:

\vspace{-0.3cm}

\begin{itemize}
	\item \texttt{\textbf{\textit{ss}}} -- масив от низове, съдържащи резултатния \textit{Graphviz} низ на части, които ще бъдат конкатенирани на по-късен етап.
	\item \texttt{\textbf{\textit{destinations}}} -- масив от низове, всеки от които представлява име на позиция
\end{itemize}

\vspace{0.55cm}

Пример за декларация на функцията \texttt{generateOutgoingEdgesFromTransition} е показан във \textbf{\autoref{lst:generate-outgoing-edges-from-transition}}.
\begin{lstlisting}[
	language=TypeScript,
	caption={Декларация на функция \textit{generateOutgoingEdgesFromTransition}.},
	label={lst:generate-outgoing-edges-from-transition}
	]
function generateOutgoingEdgesFromTransition(
    ss: string[],
    transitionName: string,
    destinations: string[],
): void
\end{lstlisting}

Функцията създава всички ребра между преход и изходните му позиции и ги добавя в резултатния масив. Не връща нищо.

Параметри на функцията:

\vspace{-0.3cm}

\begin{itemize}
	\item \texttt{\textbf{\textit{ss}}} -- масив от низове, съдържащи резултатния \textit{Graphviz} низ на части, които ще бъдат конкатенирани на по-късен етап.
	\item \texttt{\textbf{\textit{transitionName}}} -- низ, представляващ име на преход
	\item \texttt{\textbf{\textit{destionations}}} -- масив от низове, всеки от които представлява име на позиция
\end{itemize}

\vspace{0.55cm}

Пример за декларация на функцията \texttt{generateTransitionsString} е показан във \textbf{\autoref{lst:generate-transitions-string}}.
\begin{lstlisting}[
	language=TypeScript,
	caption={Декларация на функция \textit{generateTransitionsString}.},
	label={lst:generate-transitions-string}
	]
function generateTransitionsString(
    ss: string[],
    transitionToOutputs: Map<string, string[]>,
    transitionToInputs: Map<string, string[]>,
): string
\end{lstlisting}

Функцията създава всички преходи и ги добавя в резултатния масив. Връща конкатенацията на всички низовете в \textit{ss}.

Параметри на функцията:

\vspace{-0.3cm}

\begin{itemize}
	\item \texttt{\textbf{\textit{ss}}} -- масив от низове, съдържащи резултатния \textit{Graphviz} низ на части, които ще бъдат конкатенирани на по-късен етап.
	\item \texttt{\textbf{\textit{transitionToOutputs}}} -- речник с ключ име на преход и стойност всички позиции, които излизат от него.
	\item \texttt{\textbf{\textit{transitionToInput}}} -- речник с ключ име на преход и стойност всички позиции, които влизат в него.
\end{itemize}

\vspace{0.55cm}
%here i put new page command
Пример за декларация на функцията \texttt{fillMaps} е показан във \textbf{\autoref{lst:fill-maps}}.
\begin{lstlisting}[
	language=TypeScript,
	caption={Декларация на функция \textit{fillMaps}.},
	label={lst:fill-maps}
	]
function fillMaps(
    genNet: AlgorithmPlace[],
    transitionToOutputs: Map<string, string[]>,
    transitionToInputs: Map<string, string[]>,
    placeToTransition: Map<string, string>,
): void
\end{lstlisting}

Функцията попълва подадените речници на база входните данни за ОМ. Не връща нищо.

Параметри на функцията:

\vspace{-0.3cm}

\begin{itemize}
	\item \texttt{\textbf{\textit{genNet}}} -- масив от обекти от типа \textit{AlgorithmPlace}.
	\item \texttt{\textbf{\textit{transitionToOutputs}}} -- речник с ключ име на преход и стойност всички позиции, които излизат от него. Попълва се във функцията.
	\item \texttt{\textbf{\textit{transitionToInputs}}} -- речник с ключ име на преход и стойност всички позиции, които влизат в него. Попълва се във функцията.
	\item \texttt{\textbf{\textit{placeToTransition}}} -- речник с ключ име на позиция и стойност прехода, с който е свързан. Попълва се във функцията.
\end{itemize}

\vspace{0.55cm}

Пример за декларация на функцията \texttt{generateCircleNode} е показан във \textbf{\autoref{lst:generate-circle-node}}.
\begin{lstlisting}[
	language=TypeScript,
	caption={Декларация на функция \textit{generateCircleNode}.},
	label={lst:generate-circle-node}
	]
function generateCircleNode(
    res: string[],
    genNet: AlgorithmPlace[]
): void
\end{lstlisting}

Функцията създава \textit{Graphviz} низове за окръжности с определени характеристики, които да отговарят на всяка позиция. Не връща нищо.

Параметри на функцията:

\vspace{-0.3cm}

\begin{itemize}
	\item \texttt{\textbf{\textit{res}}} -- масив от низове, в който се записват \textit{Graphviz} низовете за окръжности
	\item \texttt{\textbf{\textit{genNet}}} -- масив от обекти от типа \textit{AlgorithmPlace}
\end{itemize}

\vspace{0.55cm}
%here i put new page command
Пример за декларация на функцията \texttt{getFirstFreeInvisibleNode} е показан във \textbf{\autoref{lst:get-first-free-invisible-node}}.
\begin{lstlisting}[
	language=TypeScript,
	caption={Декларация на функция \textit{getFirstFreeInvisibleNode}.},
	label={lst:get-first-free-invisible-node}
	]
function getFirstFreeInvisibleNode(
    transStr: string,
    transitionToFreeInvisNodeIndex: Map<string, number>,
): number
\end{lstlisting}

\vspace{-0.2cm}

Функцията намира номера на първата свободна невидима точка за даден преход и я маркира като използвана. Връща номер на първата свободна невидима точка.

Параметри на функцията:

\vspace{-0.3cm}

\begin{itemize}
	\item \texttt{\textbf{\textit{transStr}}} -- низ, представляващ име на преход
	\item \texttt{\textbf{\textit{transitionToFreeInvisNodeIndex}}} -- речник съдържащ първата свободна невидима точка за всеки преход
\end{itemize}

\vspace{0.55cm}

Пример за декларация на функцията \texttt{generateOutputEdgesFromPlaces} е показан във \textbf{\autoref{lst:generate-output-edges-from-places}}.
\begin{lstlisting}[
	language=TypeScript,
	caption={Декларация на функция \textit{generateOutputEdgesFromPlaces}.},
	label={lst:generate-output-edges-from-places}
	]
function generateOutputEdgesFromPlaces(
    res: string[],
    placeToTransition: Map<string, string>,
): void
\end{lstlisting}

\vspace{-0.2cm}

Функцията създава всяко ребро излизащо от позиция и влизащо в преход. Не връща нищо.

Параметри на функцията:

\vspace{-0.3cm}

\begin{itemize}
	\item \texttt{\textbf{\textit{res}}} -- масив от низове, съдържащи резултатния \textit{Graphviz} низ на части, които ще бъдат конкатенирани на по-късен етап.
	\item \texttt{\textbf{\textit{placeToTransition}}} -- речник с ключ име на позиция и стойност прехода, с който е свързан. Попълва се във функцията
\end{itemize}

\vspace{0.55cm}
%here i put new page command

Пример за декларация на функцията \texttt{isValidInput} е показан във \textbf{\autoref{lst:is-valid-input}}.
\begin{lstlisting}[
	language=TypeScript,
	caption={Декларация на функция \textit{isValidInput}.},
	label={lst:is-valid-input}
	]
function isValidInput(
    transitionToOutputs: Map<string, string[]>,
    transitionToInputs: Map<string, string[]>,
    placeToTransition: Map<string, string>,
): boolean
\end{lstlisting}

Функцията проверява дали данните в речниците представляват валидна ОМ. Връща булева стойност, отговаряща на това дали са правилни входните данни.

Параметри на функцията:

\vspace{-0.3cm}

\begin{itemize}
	\item \texttt{\textbf{\textit{transitionToOutputs}}} -- речник с ключ име на преход и стойност всички позиции, които излизат от него.
	\item \texttt{\textbf{\textit{transitionToInputs}}} -- речник с ключ име на преход и стойност всички позиции, които влизат в него.
	\item \texttt{\textbf{\textit{placeToTransition}}} -- речник с ключ име на позиция и стойност прехода, с който е свързан.
\end{itemize}

\vspace{0.55cm}

Пример за декларация на функцията \texttt{generateGraphvizString} е показан във \textbf{\autoref{lst:generate-graphviz-string}}.
\begin{lstlisting}[
	language=TypeScript,
	caption={Декларация на функция \textit{generateGraphvizString}.},
	label={lst:generate-graphviz-string}
	]
function generateGraphvizString(
    input: AlgorithmInput
): string
\end{lstlisting}

Функцията създава \textit{Graphviz} низ на ОМ на база входните данни. Връща \textit{Graphviz} низ.

Параметри на функцията:

\vspace{-0.3cm}

\begin{itemize}
	\item \texttt{\textbf{\textit{input}}} -- ОМ, представена в определен формат нужен на алгоритъма.
\end{itemize} 

\subsubsection{Разработен алгоритъм 2: преобразуване на \textit{SVG} към \textit{TeX}}

Кодът за преобразуване на \textit{SVG} низ към \textit{TeX} формат е показан във \textbf{\autoref{lst:svg-to-tex-conversion}}.

\begin{lstlisting}[
	language=TypeScript,
	caption={Алгоритъм за преобразуване на \textit{SVG} към \textit{TeX}.},
	label={lst:svg-to-tex-conversion}
	]
export function convertEdgesToXML(svgString: string) {
    const parser = new DOMParser();
    const svgDoc = parser.parseFromString(
        svgString, 
        'image/svg+xml'
    );
    const edgePaths = svgDoc.querySelectorAll('.edge path');
    
    edgePaths.forEach((pathEl) => {
        const originalD = pathEl.getAttribute('d') || '';
        // Regex to capture coordinate pairs (x,y)
        const coordRegex = /([+-]?\d+(?:\.\d+)?),?\s*/ + 
            /([+-]?\d+(?:\.\d+)?)/g;
        let match;
        const coords = [];
        
        while ((match = coordRegex.exec(originalD)) !== null) {
            const x = parseFloat(match[1]);
            const y = parseFloat(match[2]);
            coords.push([x, y]);
        }
        
        if (coords.length > 0) {
            let newD = `M ${coords[0][0]} ${coords[0][1]}`;
            for (let i = 1; i < coords.length; i++) {
                newD += ` L ${coords[i][0]} ${coords[i][1]}`;
            }
            pathEl.setAttribute('d', newD);
        }
    });
    
    // Return the modified SVG as a string
    const serializer = new XMLSerializer();
    return serializer.serializeToString(svgDoc);
}
	
/**
 * svgToTeX:
 * - Parses a Graphviz‐style SVG to extract places (circle)
 *      and transitions (rect/triangle).
 * - Forces each place circle to have 
 *      radius=0.5 in the final \LaTeX\ output.
 * - Moves place labels closer to the circle 
 *      boundary by a configurable factor.
 * - Shifts transition labels upward by a configurable offset 
 *      so they don't overlap the triangle.
 * - **Scales the entire layout by 2/3** 
 *      (so edges end up 2/3 their original length).
 * - Outputs a \LaTeX\ "picture" environment 
 *      with lines, vectors, and circles.
 *
 * @param {string} svgString
 *      - The raw SVG content as a string.
 * @param {number} [placeLabelOffsetFactor=0.85] 
 *      - Factor to scale label offset for places.
 * @param {number} [transitionLabelVerticalOffset=3.75] 
 *      - Distance to shift transition labels upward.
 * @param {number} [fixedPlaceRadius=5] 
 *      - Fixed radius for place circles in pixels.
 */
 
export function svgToTeX(
    svgString: string,
    placeLabelOffsetFactor = 0.85,
    transitionLabelVerticalOffset = 3.75,
    fixedPlaceRadius = 5, // in px
) {
    // -----------------------
    // CONFIG / CONSTANTS
    // -----------------------
    const pxToTeX = 0.1; // 1 px => 0.1 TeX units
    //const LENGTH_SCALE_FACTOR = 2 / 3; // <-- key factor: 
    //shrink geometry to 2/3
    
    // We want final circle radius=0.5 in TeX 
    // => diameter=1 => in raw SVG coords = 5 px
    // so we keep forcing that in the final output:
    //const FIXED_PLACE_RADIUS = 0.5 / pxToTeX; 
    // 0.5 / 0.1 = 5 px
    
    // Helper: parse path "M x y L x y..."
    function parsePathD(pathD: string) {
        const coordRegex = 
            /[ML]\s*([\-\d.]+)\s*,?\s*([\-\d.]+)/g;
        let match;
        const coords = [];
        while ((match = coordRegex.exec(pathD)) !== null) {
            coords.push({ 
                x: parseFloat(match[1]), 
                y: parseFloat(match[2]) 
            });
        }
        return coords;
    }
    
    // Parse the SVG as DOM
    const parser = new DOMParser();
    const svgDoc = parser.parseFromString(
        svgString, 
        'image/svg+xml'
    );
    const svgEl = svgDoc.querySelector('svg');
    if (!svgEl) {
        console.error('Not a valid SVG input.');
        return '';
    }
    
    // Possibly read global transform from g#graph0
    let globalTranslate = { x: 0, y: 0 };
    let globalScale = 1.0;
    const mainG = svgDoc.querySelector('g#graph0');
    if (mainG) {
        const tf = mainG.getAttribute('transform') || '';
        const tMatch = tf.match(/translate\(([^,]+),([^,]+)\)/);
        const sMatch = tf.match(/scale\(([^)]+)\)/);
        if (tMatch) {
            globalTranslate.x = parseFloat(tMatch[1]);
            globalTranslate.y = parseFloat(tMatch[2]);
        }
        if (sMatch) {
            globalScale = parseFloat(sMatch[1]);
        }
    }
    
    // Modified "applyGlobal" to also do the 2/3 length scaling
    function applyGlobal(x: number, y: number) {
        const X = x * globalScale + globalTranslate.x;
        const Y = y * globalScale + globalTranslate.y;
        return {
            x: X * lengthScaleFactor,
            y: Y * lengthScaleFactor,
        };
    }
    
    // -----------------------
    // 1) Node shapes
    // -----------------------
    const nodeShapes: any = {};
    
    // Places => circle
    const placeNodes = svgDoc.querySelectorAll('g.node.place');
    placeNodes.forEach((g) => {
        const titleEl = g.querySelector('title');
        const ellEl = g.querySelector('ellipse');
        if (!titleEl || !ellEl) return;
        
        const name = titleEl.textContent!.trim();
        const cx = parseFloat(ellEl.getAttribute('cx')!);
        const cy = parseFloat(ellEl.getAttribute('cy')!);
        const centerT = applyGlobal(cx, cy);
        
        nodeShapes[name] = {
            type: 'circle',
            cx: centerT.x,
            cy: centerT.y,
            r: fixedPlaceRadius,
        };
    });
    
    // Transitions => rect or triangle bounding
    const transitionNodes = svgDoc.querySelectorAll(
        'g.node.transition'
    );
    transitionNodes.forEach((g) => {
        const titleEl = g.querySelector('title');
        const polyEl = g.querySelector('polygon');
        if (!titleEl || !polyEl) return;
        
        const name = titleEl.textContent!.trim();
        const coords = polyEl.getAttribute('points')!.split(
            /\s+/
        );
        const Xs: any[] = [],
            Ys: any[] = [];
        coords.forEach((pt) => {
            const [xx, yy] = pt.split(',').map(parseFloat);
            const xyT = applyGlobal(xx, yy);
            Xs.push(xyT.x);
            Ys.push(xyT.y);
        });
        
        const xMin = Math.min(...Xs),
            xMax = Math.max(...Xs);
        const yMin = Math.min(...Ys),
            yMax = Math.max(...Ys);
        
        nodeShapes[name] = {
            type: 'rect',
            xMin,
            xMax,
            yMin,
            yMax,
        };
    });
    
    // -----------------------
    // 2) Labels
    // -----------------------
    const places: any[] = [];
    placeNodes.forEach((g) => {
        const titleEl = g.querySelector('title');
        const textEl = g.querySelector('text.node-name');
        const ellEl = g.querySelector('ellipse');
        if (!titleEl || !textEl || !ellEl) return;
        
        const pName = titleEl.textContent!.trim();
        const cx = parseFloat(ellEl.getAttribute('cx')!);
        const cy = parseFloat(ellEl.getAttribute('cy')!);
        const labelX = parseFloat(textEl.getAttribute('x')!);
        const labelY = parseFloat(textEl.getAttribute('y')!);
        
        const cT = applyGlobal(cx, cy);
        const lT = applyGlobal(labelX, labelY);
        
        // Move label slightly closer to the circle boundary
        const dx = lT.x - cT.x;
        const dy = lT.y - cT.y;
        const newLabelX = cT.x + placeLabelOffsetFactor * dx;
        const newLabelY = cT.y + placeLabelOffsetFactor * dy;
        
        places.push({
            name: pName,
            cx: cT.x,
            cy: cT.y,
            labelX: newLabelX,
            labelY: newLabelY,
        });
    });
    
    const transitions: any[] = [];
    transitionNodes.forEach((g) => {
        const titleEl = g.querySelector('title');
        const textEl = g.querySelector('text.transition-name');
        const polyEl = g.querySelector('polygon');
        if (!titleEl || !textEl || !polyEl) return;
        
        const tName = titleEl.textContent!.trim()!;
        const labelX = parseFloat(textEl.getAttribute('x')!);
        const labelY = parseFloat(textEl.getAttribute('y')!);
        const lT = applyGlobal(labelX, labelY);
        
        // For drawing vertical line + triangle
        const rawPts = polyEl.getAttribute('points')!.split(
            /\s+/
        );
        const Xs: any[] = [],
            Ys: any[] = [];
        rawPts.forEach((pt) => {
            const [xx, yy] = pt.split(',').map(parseFloat);
            const xyT = applyGlobal(xx, yy);
            Xs.push(xyT.x);
            Ys.push(xyT.y);
        });
        const xMid = (Math.min(...Xs) + Math.max(...Xs)) / 2;
        const yMin = Math.min(...Ys);
        const yMax = Math.max(...Ys);
        
        transitions.push({
            name: tName,
            // SHIFT UP 
            // so the label won't collide with the triangle
            labelX: lT.x,
            labelY: lT.y - transitionLabelVerticalOffset,
            topX: xMid,
            topY: yMin,
            botX: xMid,
            botY: yMax,
        });
    });
    
    // -----------------------
    // 3) Edges
    // -----------------------
    const edges: any[] = [];
    const edgeNodes = svgDoc.querySelectorAll('g.edge');
    edgeNodes.forEach((g) => {
        let origin = '',
            destination = '';
        
        const classStr = g.getAttribute('class') || '';
        const match = classStr.match(
            /edge\s+([\w\d]+)_{3,}([\w\d]+)/
        );
        if (match) {
            origin = match[1];
            destination = match[2];
        } else {
            // fallback: parse from <title> text "A->B"
            const titleEl = g.querySelector('title');
            if (titleEl && titleEl.textContent!.includes('->')) {
                const tStr = titleEl.textContent;
                const arrowMatch = tStr!.match(
                    /^(.*?)\s*->\s*(.*?)$/
                );
                if (arrowMatch) {
                    origin = arrowMatch[1];
                    destination = arrowMatch[2];
                }
            }
        }
        
        const pathEl = g.querySelector('path');
        if (!pathEl) return;
        const dAttr = pathEl.getAttribute('d') || '';
        const rawPts = parsePathD(dAttr).map(
            (pt) => applyGlobal(pt.x, pt.y)
        );
        
        // arrow polygon
        let arrowPoly: any[] = [];
        const polyEl = g.querySelector('polygon');
        if (polyEl) {
            const polyPts = polyEl.getAttribute('points')!.split(
                /\s+/
            );
            arrowPoly = polyPts.map((pt) => {
                const [xx, yy] = pt.split(',').map(parseFloat);
                return applyGlobal(xx, yy);
            });
        }
        
        edges.push({
            origin,
            destination,
            rawPoints: rawPts,
            arrowPoly,
        });
    });
    
    // -----------------------
    // 4) Clip edges to node boundary
    // -----------------------
    function fixCircleBoundary(
        cx: number, 
        cy: number, 
        r: number, 
        oldPt: any
    ) {
        const dx = oldPt.x - cx;
        const dy = oldPt.y - cy;
        let dist = Math.sqrt(dx * dx + dy * dy);
        if (dist < 1e-9) dist = 1;
        const scale = r / dist;
        return { x: cx + scale * dx, y: cy + scale * dy };
    }
    
    function fixRectBoundary(
        xMin: number,
        xMax: number,
        yMin: number,
        yMax: number,
        oldPt: any,
        neighborPt: any,
    ) {
        const dx = neighborPt.x - oldPt.x;
        const dy = neighborPt.y - oldPt.y;
        // Decide whether to clamp horizontally or vertically
        if (Math.abs(dx) >= Math.abs(dy)) {
            // horizontal
            if (dx > 0) return { x: xMax, y: oldPt.y };
            else return { x: xMin, y: oldPt.y };
        } else {
            // vertical
            if (dy > 0) return { x: oldPt.x, y: yMax };
            else return { x: oldPt.x, y: yMin };
        }
    }
    
    function fixEdgeEndpoints(e: any) {
        const { origin, destination, rawPoints, arrowPoly } = e;
        if (rawPoints.length < 2) return;
        
        // fix first
        const shapeA = nodeShapes[origin];
        if (shapeA) {
            const p0 = rawPoints[0];
            const p1 = rawPoints[1];
            if (shapeA.type === 'circle') {
                rawPoints[0] = fixCircleBoundary(
                    shapeA.cx, shapeA.cy, shapeA.r, p0
                );
            } else {
                rawPoints[0] = fixRectBoundary(
                    shapeA.xMin,
                    shapeA.xMax,
                    shapeA.yMin,
                    shapeA.yMax,
                    p0,
                    p1,
                );
            }
        }
        
        // fix last
        const shapeB = nodeShapes[destination];
        if (shapeB) {
            const n = rawPoints.length;
            const pLast = rawPoints[n - 1];
            const pBefore = rawPoints[n - 2];
            let newLast = pLast;
            if (shapeB.type === 'circle') {
                newLast = fixCircleBoundary(
                    shapeB.cx, shapeB.cy, shapeB.r, pLast
                );
            } else {
                newLast = fixRectBoundary(
                    shapeB.xMin,
                    shapeB.xMax,
                    shapeB.yMin,
                    shapeB.yMax,
                    pLast,
                    pBefore,
                );
            }
            rawPoints[n - 1] = newLast;
            
            // shift arrow polygon 
            // so its tip is at the new endpoint
            if (arrowPoly && arrowPoly.length > 0) {
                const oldTip = arrowPoly[0];
                const dx = newLast.x - oldTip.x;
                const dy = newLast.y - oldTip.y;
                for (let i = 0; i < arrowPoly.length; i++) {
                    arrowPoly[i] = {
                        x: arrowPoly[i].x + dx,
                        y: arrowPoly[i].y + dy,
                    };
                }
            }
        }
    }
    
    edges.forEach((e) => fixEdgeEndpoints(e));
    
    // -----------------------
    // 5) Compute bounding box
    // -----------------------
    const allX = [],
        allY = [];
    places.forEach((p) => {
        allX.push(p.cx, p.labelX);
        allY.push(p.cy, p.labelY);
    });
    transitions.forEach((t) => {
        allX.push(t.labelX, t.topX, t.botX);
        allY.push(t.labelY, t.topY, t.botY);
    });
    edges.forEach((e) => {
        e.rawPoints.forEach((pt: any) => {
            allX.push(pt.x);
            allY.push(pt.y);
        });
        e.arrowPoly.forEach((pt: any) => {
            allX.push(pt.x);
            allY.push(pt.y);
        });
    });
    
    if (allX.length === 0) {
        allX.push(0, 100);
        allY.push(0, 100);
    }
    const minAllX = Math.min(...allX);
    const maxAllX = Math.max(...allX);
    const minAllY = Math.min(...allY);
    const maxAllY = Math.max(...allY);
    
    const rangeX = maxAllX - minAllX;
    const rangeY = maxAllY - minAllY;
    
    // Convert to TeX coords (0,0 at lower-left corner):
    function toTeXCoords(x: number, y: number) {
        const tx = (x - minAllX) * pxToTeX;
        const ty = (maxAllY - y) * pxToTeX; // flip Y
        return { x: tx, y: ty };
    }
    
    const picW = rangeX * pxToTeX;
    const picH = rangeY * pxToTeX;
    
    // -----------------------
    // 6) Build lines in TeX
    // -----------------------
    function segmentsToTeX(points: any) {
        const lines = [];
        for (let i = 0; i < points.length - 1; i++) {
            const p1 = points[i];
            const p2 = points[i + 1];
            const isLast = i === points.length - 2;
            
            const dx = p2.x - p1.x;
            const dy = p2.y - p1.y;
            const cmd = isLast ? '\\vector' : '\\line';
            
            const prefix = 
                `\\put(${p1.x.toFixed(2)},${p1.y.toFixed(2)})`;
            if (Math.abs(dx) > Math.abs(dy)) {
                // horizontal
                const L = Math.abs(dx).toFixed(2);
                const dir = dx >= 0 ? 1 : -1;
                lines.push(`${prefix}{${cmd}(${dir},0){${L}}}`);
            } else {
                // vertical
                const L = Math.abs(dy).toFixed(2);
                const dir = dy >= 0 ? 1 : -1;
                lines.push(`${prefix}{${cmd}(0,${dir}){${L}}}`);
            }
        }
        return lines.join('\n');
    }
    
    // -----------------------
    // 7) Final TeX
    // -----------------------
    const texLines = [];
    texLines.push('\\tiny');
    texLines.push('\\begin{center}');
    texLines.push('\\setlength{\\unitlength}{8pt}');
    texLines.push(
        `\\begin{picture}(${picW.toFixed(2)},${picH.toFixed(2)})`
    );
    texLines.push('\\linethickness{0.4pt}');
    
    // (a) Transitions
    texLines.push('% --- Transitions ---');
    transitions.forEach((t) => {
        const labelPos = toTeXCoords(t.labelX, t.labelY);
        const topPos = toTeXCoords(t.topX, t.topY);
        const botPos = toTeXCoords(t.botX, t.botY);
        
        // Label (anchored bottom so it extends upward)
        texLines.push(
            `\\put(${labelPos.x.toFixed(2)},` +
            `${labelPos.y.toFixed(2)})` +
            `{\\makebox(0,0)[b]{${t.name}}}`
        );
        // Triangle
        texLines.push(
            `\\put(${topPos.x.toFixed(2)},` + 
            `${topPos.y.toFixed(2)})` +
            `{\\makebox(0,0)[bb]{\\Large$\\bigtriangledown$}}`
        );
        // Vertical line
        const lineLen = (topPos.y - botPos.y).toFixed(2);
        texLines.push(
            `\\put(${topPos.x.toFixed(2)},` + 
            `${topPos.y.toFixed(2)})` +
            `{\\line(0,-1){${lineLen}}}`
        );
    });
    
    // (b) Places
    texLines.push('% --- Places ---');
    places.forEach((p) => {
        const cPos = toTeXCoords(p.cx, p.cy);
        const lPos = toTeXCoords(p.labelX, p.labelY);

        // label
        texLines.push(
            `\\put(${lPos.x.toFixed(2)},${lPos.y.toFixed(2)})` + 
            `{\\makebox(0,0)[cc]{${p.name}}}`
        );
        // circle with diameter=1 => radius=0.5
        texLines.push(
            `\\put(${cPos.x.toFixed(2)},${cPos.y.toFixed(2)})` + 
            `{\\circle{1.0}}`
        );
    });
    
    // (c) Edges
    texLines.push('% --- Edges ---');
    edges.forEach((e) => {
        const { origin, destination, rawPoints } = e;
        if (!rawPoints || rawPoints.length < 2) return;
        
        texLines.push(`%% ${origin} -> ${destination}`);
        const Tpoints = rawPoints.map(
            (pt: any) => toTeXCoords(pt.x, pt.y)
        );
        texLines.push(segmentsToTeX(Tpoints));
    });
    
    texLines.push('\\end{picture}');
    texLines.push('\\end{center}');
    texLines.push('\\normalsize');
    
    return texLines.join('\n');
}

// ------------------------------------------------------
// EXAMPLE USAGE:
//
// const exampleSvg = `... your SVG content ...`;
// const fixedSvg = convertEdgesToXML(exampleSvg);
// const texOutput = svgToTeX(fixedSvg, 0.80, 4.0);
//
// console.log(texOutput);
\end{lstlisting}

%here i put new page command
\paragraph{Декларации на функциите.}
\mbox{}

Пример за декларация на функцията \texttt{convertEdgesToXML} е показан във \textbf{\autoref{lst:convert-edges-to-xml}}.
\begin{lstlisting}[
	language=TypeScript,
	caption={Декларация на функция \textit{convertEdgesToXML}.},
	label={lst:convert-edges-to-xml}
	]
function convertEdgesToXML(svgString: string)
\end{lstlisting}

Функцията получава \textit{SVG} като низ и рефакторира част от ребрата в него, така че да бъдат по-удобни за обработка от \textit{svgToTex} функцията.

\vspace{0.55cm}

Пример за декларация на функцията \texttt{svgToTeX} е показан във \textbf{\autoref{lst:svg-to-tex}}.
\begin{lstlisting}[
	language=TypeScript,
	caption={Декларация на функция \textit{svgToTeX} за преобразуване на \textit{SVG} в \textit{TeX}.},
	label={lst:svg-to-tex}
	]
function svgToTeX(
    svgString: string,
    placeLabelOffsetFactor = 0.85,
    transitionLabelVerticalOffset = 3.75,
    fixedPlaceRadius = 5, // in px
)
\end{lstlisting}

Функцията преобразува \textit{SVG} низ, представляващ визуализация на ОМ към \textit{TeX} формат. 

Параметри на функцията:

\vspace{-0.3cm}

\begin{itemize}
	\item \texttt{\textbf{\textit{svgString}}} -- \textit{SVG} низ представляващ визуализация на ОМ.
	\item \texttt{\textbf{\textit{placeLabelOffsetFactor}}} -- Число, определящо отстоянието между позицията и нейния етикет. Стойността по подразбиране е 0.85 пиксела.
	\item \texttt{\textbf{\textit{transitionLabelVerticalOffset}}} -- Число, представляващо коефициента на отстоянието от прехода на етикета, който го описва. Стойността по подразбиране е 3.75 пиксела.
	\item \texttt{\textbf{\textit{lengthScaleFactor}}} -- Число, представляващо коефициента на мащабиране на дължините в \textit{SVG}. Стойността по подразбиране е 0.5 пиксела.
	\item \texttt{\textbf{\textit{fixedPlaceRadius}}} -- Число в пиксели, представляващо фиксиран радиус на кръговете, представляващи позициите. Стойността по подразбиране е 5 пиксела.
\end{itemize}

%here i put new page command
\newpage

\subsection{Интерфейс и потребителско изживяване}

\hspace*{2.5em}Настоящата глава представя реализацията на \ref{task:web-ui} и \ref{task:predicates-ui}. Разглежда се дизайнът на потребителския интерфейс и основните принципи на потребителското изживяване при работа със системата \textit{OnlineGN}. Акцентът е поставен върху навигацията, визуалните компоненти и начина, по който интерфейсът подпомага интерактивната работа със симулатора, включително механизмите за визуализация, персонализация и редакция на специфичните ОМ параметри.

\subsubsection{Навигация и визуални компоненти на уеб клиента}

\hspace*{2.5em}Следващото изложение представя детайлно решението на \ref{task:web-ui}. Основният интерфейс на клиентското приложение събира всички основни функционалности на системата в една страница, както се вижда на \textbf{Фигура~\ref{fig:ui-main-screen}}.
\begin{figure}[H]
	\centering
	\setlength{\fboxrule}{2pt}
	\setlength{\fboxsep}{2pt}  
	\fbox{\includegraphics[width=1\linewidth]{"3. onlineGn/images/ui-main-screen.png"}}
	\caption{Основно работно поле и контроли.}
	\label{fig:ui-main-screen}
\end{figure}


Основните действия върху мрежата се извършват чрез функционалните бутони, илюстрирани на \textbf{Фигура~\ref{fig:ui-action-buttons}}.
\begin{figure}[H]
	\centering
	\setlength{\fboxrule}{2pt}
	\setlength{\fboxsep}{0pt}
	\fbox{\includegraphics[width=0.8\linewidth]{"3. onlineGn/images/ui-action-buttons.png"}}
	\caption{Функционални бутони в клиентския интерфейс.}
	\label{fig:ui-action-buttons}
\end{figure}

\newpage
Визуализацията на мрежата и текущата стъпка на симулацията е показана на \textbf{Фигура~\ref{fig:ui-network-and-step}}.
\begin{figure}[H]
	\centering
	\setlength{\fboxrule}{2pt}
	\setlength{\fboxsep}{3pt}
	\fbox{\includegraphics[width=1\linewidth]{"3. onlineGn/images/ui-network-and-step.png"}}
    \captionsetup{width=0.9\textwidth}
	\caption{Визуализация на мрежата и индикатор за текуща стъпка; числата в позициите показват броя ядра.}
	\label{fig:ui-network-and-step}
\end{figure}

\vspace{0.3cm}

Потребителският код се въвежда и изпълнява чрез специализиран панел, показан на \textbf{Фигура~\ref{fig:ui-user-code-panel}}.
\begin{figure}[H]
	\centering
	\setlength{\fboxrule}{2pt}
	\setlength{\fboxsep}{0pt}
	\fbox{\includegraphics[width=1\linewidth]{"3. onlineGn/images/ui-user-code-panel.png"}}
	\caption{Блок на потребителски \textit{JavaScript} код (компилиране/изпълнение).}
	\label{fig:ui-user-code-panel}
\end{figure}

\newpage
Изпълнението на единична стъпка от симулацията се инициира чрез бутона „\textit{Make a step}“, показан на \textbf{Фигура~\ref{fig:ui-make-step}}.
\begin{figure}[H]
	\centering
	\setlength{\fboxrule}{2pt}
	\setlength{\fboxsep}{3pt}
	\fbox{\includegraphics[width=1\linewidth]{"3. onlineGn/images/ui-make-step.png"}}
	\captionsetup{width=0.8\textwidth}
	\caption{Изпълнение на единична стъпка \textit{(„Make a step“)}; увеличаване на брояча.}
	\label{fig:ui-make-step}
\end{figure}

\vspace{0.3cm}
Резултатът от изпълнението на стъпка е илюстриран на \textbf{Фигура~\ref{fig:ui-after-step}}.
\begin{figure}[H]
	\centering
	\setlength{\fboxrule}{2pt}
	\setlength{\fboxsep}{3pt}
	\fbox{\includegraphics[width=1\linewidth]{"3. onlineGn/images/ui-after-step.png"}}
	\caption{Състояние на мрежата след изпълнение на стъпка.}
	\label{fig:ui-after-step}
\end{figure}

\newpage
Следващото изложение представя решението на \ref{task:predicates-ui}. Режимът на редактиране на координати се активира, както е показано на \textbf{Фигура~\ref{fig:ui-edit-positions-mode}}.

\begin{figure}[H]
	\centering
	\setlength{\fboxrule}{2pt}
	\setlength{\fboxsep}{3pt}
	\fbox{\includegraphics[width=1\linewidth]{"3. onlineGn/images/ui-edit-positions-mode.png"}}
	\caption{Режим за промяна на координати \textit{(„Edit positions“)}.}
	\label{fig:ui-edit-positions-mode}
\end{figure}

\vspace{0.3cm}

Селекцията на позиция за редактиране е показана на \textbf{Фигура~\ref{fig:ui-select-place}} (избрана е позиция \texttt{p6}).
\begin{figure}[H]
	\centering
	\setlength{\fboxrule}{2pt}
	\setlength{\fboxsep}{3pt}
	\fbox{\includegraphics[width=1\linewidth]{"3. onlineGn/images/ui-select-place.png"}}
	\caption{Селекция на позиция за редактиране (избрана е позиция p6).}
	\label{fig:ui-select-place}
\end{figure}

\newpage
Ефектът от промяната на координата на позиция е илюстриран на \textbf{Фигура~\ref{fig:ui-place-y-updated}}.

\vspace{-0.2cm}

\begin{figure}[H]
	\centering
	\setlength{\fboxrule}{2pt}
	\setlength{\fboxsep}{3pt}
	\fbox{\includegraphics[width=0.85\linewidth]{"3. onlineGn/images/ui-place-y-updated.png"}}
	\caption{Резултат след промяна на y-координатата на позиция.}
	\label{fig:ui-place-y-updated}
\end{figure}

Редактирането на преход се инициира чрез селекция, показана на \textbf{Фигура~\ref{fig:ui-select-transition}} (избран е преход T2).

\vspace{-0.2cm}

\begin{figure}[H]
	\centering
	\setlength{\fboxrule}{2pt}
	\setlength{\fboxsep}{3pt}
	\fbox{\includegraphics[width=0.85\linewidth]{"3. onlineGn/images/ui-select-transition.png"}}
	\caption{Селекция на преход за редактиране (избран е преход T2).}
	\label{fig:ui-select-transition}
\end{figure}

Резултатът от промяната на параметрите на преход е показан на \textbf{Фигура~\ref{fig:ui-transition-updated}}.

\vspace{-0.2cm}

\begin{figure}[H]
	\centering
	\setlength{\fboxrule}{2pt}
	\setlength{\fboxsep}{3pt}
	\fbox{\includegraphics[width=0.85\linewidth]{"3. onlineGn/images/ui-transition-updated.png"}}
	\caption{Резултат след промяна на дължина и x-координатата на преход.}
	\label{fig:ui-transition-updated}
\end{figure}


Ребрата могат да бъдат селектирани за редактиране, както е показано на \textbf{Фигура~\ref{fig:ui-select-edge}}.

\begin{figure}[H]
	\centering
	\setlength{\fboxrule}{2pt}
	\setlength{\fboxsep}{3pt}
	\fbox{\includegraphics[width=1\linewidth]{"3. onlineGn/images/ui-select-edge.png"}}
	\caption{Селекция на ребро за редактиране.}
	\label{fig:ui-select-edge}
\end{figure}

\vspace{0.3cm}

Добавянето и редактирането на контролни точки на ребро са илюстрирани на \textbf{Фигура~\ref{fig:ui-edge-add-point}}.

\begin{figure}[H]
	\centering
	\setlength{\fboxrule}{2pt}
	\setlength{\fboxsep}{3pt}
	\fbox{\includegraphics[width=1\linewidth]{"3. onlineGn/images/ui-edge-add-point.png"}}
	\caption{Добавяне на нова точка и промяна на координати на точки 1 и 2.}
	\label{fig:ui-edge-add-point}
\end{figure}

\newpage
Разместването на контролни точки води до промяна във формата на реброто, което е показано на \textbf{Фигура~\ref{fig:ui-edge-swap-points}}.

\begin{figure}[H]
	\centering
	\setlength{\fboxrule}{2pt}
	\setlength{\fboxsep}{3pt}
	\fbox{\includegraphics[width=1\linewidth]{"3. onlineGn/images/ui-edge-swap-points.png"}}
	\caption{Резултат след разместване на точки 1 и 2.}
	\label{fig:ui-edge-swap-points}
\end{figure}

\vspace{0.3cm}

Изтриването на контролна точка е показано на \textbf{Фигура~\ref{fig:ui-edge-delete-point}}.

\begin{figure}[H]
	\centering
	\setlength{\fboxrule}{2pt}
	\setlength{\fboxsep}{3pt}
	\fbox{\includegraphics[width=1\linewidth]{"3. onlineGn/images/ui-edge-delete-point.png"}}
	\caption{Резултат при изтриване на точка 1.}
	\label{fig:ui-edge-delete-point}
\end{figure}

\newpage
Предикатите на преход се визуализират контекстно при посочване с мишка (\textbf{Фигура~\ref{fig:ui-transition-predicate-hover}}).


\begin{figure}[H]
	\centering
	\setlength{\fboxrule}{2pt}
	\setlength{\fboxsep}{3pt}
	\fbox{\includegraphics[width=0.95\linewidth]{"3. onlineGn/images/ui-transition-predicate-hover.png"}}
	\caption{Показване на предикати на преход при посочване с мишка.}
	\label{fig:ui-transition-predicate-hover}
\end{figure}

Всички ядра и техните характеристики се показват в отделен диалогов прозорец (\textbf{Фигура~\ref{fig:ui-all-tokens-dialog}}).
\begin{figure}[H]
	\centering
	\setlength{\fboxrule}{2pt}
	\setlength{\fboxsep}{2pt}
	\fbox{\includegraphics[width=0.5\linewidth]{"3. onlineGn/images/ui-all-tokens-dialog.png"}}
	\caption{Прозорец с всички ядра и техните характеристики.}
	\label{fig:ui-all-tokens-dialog}
\end{figure}

\newpage
Настройките на симулацията се конфигурират чрез специализиран диалогов прозорец (\textbf{Фигура~\ref{fig:ui-settings-dialog}}).
\vspace{-0.3cm}
\begin{figure}[H]
	\centering
	\setlength{\fboxrule}{2pt}
	\setlength{\fboxsep}{0pt}
	\fbox{\includegraphics[width=0.5\linewidth]{"3. onlineGn/images/ui-settings-dialog.png"}}
	\caption{Настройки на симулацията; промяна на скоростта.}
	\label{fig:ui-settings-dialog}
\end{figure}

\vspace{-0.2cm}
Прозорец, показващ регистър на всяка стъпка на алгоритъма -- показва как е приоритизирал преместванията на ядра спрямо правилата и капацитетите (показва се чрез натискане на бутона \textit{„Logs“}). Прозорецът след изпълнението на една стъпка от симулацията е показан на \textbf{Фигура~\ref{fig:ui-logs-dialog}}.
\vspace{-0.3cm}
\begin{figure}[H]
	\centering
	\setlength{\fboxrule}{2pt}
	\setlength{\fboxsep}{0pt}
	\fbox{\includegraphics[width=0.5\linewidth]{"3. onlineGn/images/ui-logs-dialog.png"}}\\
	\captionsetup{width=0.8\textwidth}
	\caption{Регистър на стъпка: приоритизация на премествания спрямо правила и капацитет.}
	\label{fig:ui-logs-dialog}
\end{figure}


При неуспешно съхранение системата визуализира съобщение за грешка (\textbf{Фигура~\ref{fig:ui-save-error}}).
\begin{figure}[H]
	\centering
	\setlength{\fboxrule}{2pt}
	\setlength{\fboxsep}{0pt}
	\fbox{
\includegraphics[width=0.95\linewidth]{"3. onlineGn/images/ui-save-error.png"}}
	\captionsetup{width=0.8\textwidth}
	\caption{Неуспешен опит за съхранение на мрежа в базата от данни (пример за грешка).}
	\label{fig:ui-save-error}
\end{figure}

\vspace{0.3cm}

Успешното съхранение на мрежа генерира хипервръзка, показана на \textbf{Фигура~\ref{fig:ui-save-success}}.
\begin{figure}[H]
	\centering
	\setlength{\fboxrule}{2pt}
	\setlength{\fboxsep}{0pt}
	\fbox{\includegraphics[width=0.95\linewidth]{"3. onlineGn/images/ui-save-success.png"}}
	\captionsetup{width=0.8\textwidth}
	\caption{Успешно запазване на мрежа; генерирана хипервръзка на запазената мрежа.}
	\label{fig:ui-save-success}
\end{figure}

\newpage
Експортирането в \textit{TeX} формат се извършва чрез диалоговия прозорец на \textbf{Фигура~\ref{fig:ui-export-tex-dialog}}.
\begin{figure}[H]
	\centering
	\setlength{\fboxrule}{2pt}
	\setlength{\fboxsep}{0pt}
	\fbox{\includegraphics[width=0.5\linewidth]{"3. onlineGn/images/ui-export-tex-dialog.png"}}
	\caption{Експортиране в \textit{TeX} формат.}
	\label{fig:ui-export-tex-dialog}
\end{figure}

\vspace{0.3cm}

Генерираният \textit{TeX} резултат е показан на \textbf{Фигура~\ref{fig:ui-export-tex-result}}.
\begin{figure}[H]
	\centering
	\setlength{\fboxrule}{2pt}
	\setlength{\fboxsep}{0pt}
	\fbox{\includegraphics[width=0.5\linewidth]{"3. onlineGn/images/ui-export-tex-result.png"}}
	\caption{Генериран \textit{TeX} след избор на позволени характеристики.}
	\label{fig:ui-export-tex-result}
\end{figure}

\newpage
Мрежата може да бъде експортирана във файл, както е показано на \textbf{Фигура~\ref{fig:ui-export-file}}.
\begin{figure}[H]
	\centering
	\setlength{\fboxrule}{2pt}
	\setlength{\fboxsep}{0pt}
	\fbox{
\includegraphics[width=0.8\linewidth]{"3. onlineGn/images/ui-export-file.png"}}
	\caption{Експортиране във файл.}
	\label{fig:ui-export-file}
\end{figure}

Зареждането на предварително запазена мрежа е илюстрирано на \textbf{Фигура~\ref{fig:ui-load-from-file}}.
\begin{figure}[H]
	\centering
	\setlength{\fboxrule}{2pt}
	\setlength{\fboxsep}{3pt}
	\fbox{\includegraphics[width=0.9\linewidth]{"3. onlineGn/images/ui-load-from-file.png"}}
	\caption{Зареждане на предварително запазена мрежа (пример: \textit{export (8).json}).}
	\label{fig:ui-load-from-file}
\end{figure}

Грешка при компилиране на потребителски код е показана на \textbf{Фигура~\ref{fig:ui-compile-error}}.
\begin{figure}[H]
	\centering
	\setlength{\fboxrule}{2pt}
	\setlength{\fboxsep}{3pt}
	\fbox{\includegraphics[width=0.85\linewidth]{"3. onlineGn/images/ui-compile-error.png"}}
	\caption{Грешка при компилиране на потребителския код след промяна.}
	\label{fig:ui-compile-error}
\end{figure}

Успешното добавяне на функция в кода е илюстрирано на \textbf{Фигура~\ref{fig:ui-compile-success-add-func}}.
\begin{figure}[H]
	\centering
	\setlength{\fboxrule}{2pt}
	\setlength{\fboxsep}{3pt}
	\fbox{\includegraphics[width=0.85\linewidth]{"3. onlineGn/images/ui-compile-success-add-func.png"}}
	\caption{Успешно добавяне на нова функция в компилирания код (пример: test2).}
	\label{fig:ui-compile-success-add-func}
\end{figure}

\vspace{-0.2cm}

\textbf{Фигура~\ref{fig:ui-edit-transition-dialog}} илюстрира диалоговия прозорец за редактиране на преход, чрез който се конфигурират следните параметри:
\vspace{-0.2cm}
\begin{itemize}
	\item функция за разделяне на ядрото;
	\item приоритет на прехода;
	\item функция на преход между две позиции;
	\item капацитет на прехода между две позиции.
\end{itemize}

\begin{figure}[H]
	\centering
	\setlength{\fboxrule}{2pt}
	\setlength{\fboxsep}{0pt}
	\fbox{\includegraphics[width=0.45\linewidth]{"3. onlineGn/images/ui-edit-transition-dialog.png"}}
	\caption{Диалогов прозорец за промяна на преход.}
	\label{fig:ui-edit-transition-dialog}
\end{figure}


Предикатната матрица на преход е визуализирана на \textbf{Фигура~\ref{fig:ui-transition-matrix}}.
\begin{figure}[H]
	\centering
	\setlength{\fboxrule}{2pt}
	\setlength{\fboxsep}{0pt}
	\fbox{\includegraphics[width=0.55\linewidth]{"3. onlineGn/images/ui-transition-matrix.png"}}
	\caption{Предикатна матрица на конкретен преход.}
	\label{fig:ui-transition-matrix}
\end{figure}

\vspace{0.3cm}

Параметрите на позиция могат да бъдат редактирани чрез диалоговия прозорец, показан на \textbf{Фигура~\ref{fig:ui-edit-place-dialog}}, като се задават:

\vspace{-0.3cm}

\begin{itemize}
	\item функция на сливане;
	\item капацитет на позицията;
	\item приоритет на позицията;
	\item характеристична функция.
\end{itemize}

\begin{figure}[H]
	\centering
	\setlength{\fboxrule}{2pt}
	\setlength{\fboxsep}{0pt}
	\fbox{\includegraphics[width=0.55\linewidth]{"3. onlineGn/images/ui-edit-place-dialog.png"}}
	\caption{Прозорец за промяна на позиция.}
	\label{fig:ui-edit-place-dialog}
\end{figure}

Допълнителни подробности за позицията са представени на \textbf{Фигура~\ref{fig:ui-place-details}}.

\begin{figure}[H]
	\centering
	\setlength{\fboxrule}{2pt}
	\setlength{\fboxsep}{0pt}
	\fbox{\includegraphics[width=0.55\linewidth]{"3. onlineGn/images/ui-place-details.png"}}
	\captionsetup{width=0.8\textwidth}
	\caption{Детайли за позиция – визуализация на параметри/ядра (допълнение към предходния диалогов прозорец).}
	\label{fig:ui-place-details}
\end{figure}



\subsubsection{Помощна система, документация и локализация}
\hspace*{2.5em}С цел улесняване на потребителите при работа със системата е предвидена помощна система, допълваща функционалността на уеб клиента. Тя включва вградена документация, контекстна помощ и възможности за локализация на интерфейса, което позволява по-широка приложимост на симулатора в образователен и изследователски контекст.

\subsection{Тестове и валидация}
\hspace*{2.5em}Тестовете и валидацията са ключови елементи за осигуряване на качеството, надеждността и възпроизводимостта на разработения \textbf{симулатор за ОМ}. Спецификата на областта изисква коректност не само на отделни програмни компоненти, но и на \textbf{семантиката на симулационния процес} (например изпълнение на преходи, придвижване/разделяне/сливане на ядра, актуализация на състояния и атрибути, спазване на правилата за допустимост на преходите и детерминистичен ред на обработка при еднакви входни условия). Поради това е приложена многостепенна тестова стратегия, която комбинира автоматизирани тестове за регресия и повторяемост, интеграционни проверки на критичните интерфейси и независима ръчна валидация на потребителско ниво.

\begin{itemize}
    \item \textbf{Тестове на ниво програмна единица (\textit{Unit tests})} -- покриват функции, класове и модули, отговорни за: валидиране на входни данни за ОМ; операции върху структурата на мрежата; изпълнение на характеристични фунцкии; алгоритми (например за движение на ядрата). За реализацията на unit тестовете в среда \textit{TypeScript/Node.js} е използван \textit{Jest} (с \textit{TypeScript} интеграция), а измерването на тестово покритие се извършва чрез \textit{Istanbul/nyc} (\textit{coverage} репорти в \textit{CI}). \textit{Unit} тестовете се изпълняват автоматично в \textit{CI pipeline} при всяка промяна по кода с цел ранно откриване на регресии; поддържат се прагове за минимално допустимо покритие (на ниво линии/функции/клонове), така че при бъдещи изменения да не се допуска неконтролирано понижаване на тестовия обхват.

    \item \textbf{Интеграционно тестване (\textit{Integration testing})} -- валидира взаимодействието между основните подсистеми на симулатора и устойчивостта на изпълнение при реалистични конфигурации. Интеграционните тестове са фокусирани върху:
    \begin{itemize}
        \item \textbf{сървър–база данни} (устойчиво съхранение и извличане на модели на ОМ и симулационни резултати, коректност на сериализация/десериализация, транзакционно поведение и обработка на грешки);
        \item \textbf{клиент–сървър} (коректност на \textit{API} договорите, валидиране на входовете, кодове за статус, поведение при частични откази/таймаути и устойчивост на потребителските потоци при типични сценарии).
    \end{itemize}

    \item \textbf{Автоматизирано тестване от край до край (\textit{Automated E2E testing})} -- проверява критични потребителски потоци в \textit{SPA} интерфейса в контролирана среда, като симулира действия на потребител (зареждане на ОМ, стартиране/управление на симулация, визуализация на резултати, обработка на невалидни входове). \textit{E2E} тестовете се изпълняват като финален автоматизиран слой преди доставка на версия; използван е инструментариум от тип \textit{Playwright} за \textit{UI} автоматизация, което осигурява повторяемост и проследимост на резултатите.

    \item \textbf{Независима ръчна валидация (\textit{Manual acceptance/E2E validation})} -- ръчните проверки с цел \emph{валидация} са проведени от \textbf{независими оценители}, които не участват пряко в имплементацията на симулатора. Това е необходимо, тъй като самопроверката от автора/разработчика не представлява независима оценка и създава риск от пропускане на дефекти поради предварителни очаквания към поведението на системата. На оценителите е предоставен тестов протокол (контролен списък) със сценарии за приемане и критерии за коректно поведение; установените несъответствия са документирани като дефекти с конкретни стъпки за възпроизвеждане и очакван/реален резултат. Проведената независима валидация потвърждава коректното изпълнение на основните потребителски сценарии и подпомага отстраняването на установени несъответствия в хода на разработката. Допълнително, по време на разработка се изпълняват кратки \textit{smoke} проверки от разработчика след промени, които се разглеждат като оперативна самопроверка, а не като формална валидация.
\end{itemize}

\textbf{Тестови модели: примерни ОМ с референтно поведение.}

За да бъде тестването на симулационната семантика систематично, са разработени \textbf{примерни ОМ} (от \cite{Koycheva2002}), предназначени специално за тестови и валидационни цели. За всяка от тях са дефинирани очаквани резултати, получени по един от следните подходи:
\begin{itemize}
    \item \textbf{Теоретично проиграване} (ръчно/формално проследяване на стъпките на симулация при зададени начални условия и правила), което служи като „еталон“ за коректно поведение;
    \item \textbf{Крос--валидация} чрез изпълнение на същите мрежи в \textbf{други налични симулатори за ОМ} и сравнение на ключови резултати (например активирани преходи по стъпки, крайно разпределение на ядра, промени в атрибути и състояния), при еквивалентни начални параметри.
\end{itemize}


Част от примерните ОМ са проиграни както теоретично, така и чрез крос--валидация в други налични симулационни среди, което дава допълнителна увереност в коректността на получените резултати. Сравнителният анализ показва пълно съвпадение по отношение на основните стъпки на изпълнение и крайните състояния при еднакви начални условия.

Тестовият набор включва както \textbf{референтни мрежи} с предвидим резултат (подходящи за регресионни проверки), така и \textbf{гранични и негативни случаи}: минимални/максимални конфигурации, празни или силно разклонени структури, невалидни входове (например липсващи задължителни елементи, неконсистентни референции, противоречиви условия). По този начин се осигурява проверка както на алгоритмичната коректност на симулационното ядро, така и на устойчивото поведение на системата при нетипични и грешни входни данни.

Комбинацията от автоматизирани \textit{unit/integration/E2E} тестове, независима ръчна валидация и използването на примерни ОМ с референтно поведение предоставя проследим и възпроизводим механизъм за оценка на качеството на разработения симулатор и за ограничаване на регресии при бъдещо развитие на системата.

\paragraph{Оценка на степента на постигане на нефункционалните изисквания.}
\mbox{} \\
Посочените в \textbf{раздел~\ref{sec:non-functional-requirements}} нефункционални изисквания са взети предвид както при проектирането на архитектурата, така и при реализирането на основните модули на системата. В рамките на разработения симулатор те са постигнати на ниво, съответстващо на предвидените сценарии на работа и характера на приложението. По-конкретно, архитектурното разделяне на отговорностите, използването на утвърдени технологии и реализираните механизми за визуализация, симулация, съхраняване и обмен на данни дават основание да се приеме, че системата удовлетворява поставените изисквания за производителност, надеждност, сигурност, използваемост и разширяемост.